%==============================================================
% CSE 3200: System Development Project
% Report - Invigilo (Secure Exam Desktop Application)
% Khulna University of Engineering & Technology
%==============================================================

\documentclass[12pt, a4paper]{report}

%--------------------------------------------------------------
% PACKAGES
%--------------------------------------------------------------
\usepackage[T1]{fontenc}
\usepackage{mathptmx}           % Times Roman font
\usepackage[utf8]{inputenc}
\usepackage{geometry}
\usepackage{setspace}
\usepackage{titlesec}
\usepackage{tocloft}
\usepackage{graphicx}
\usepackage{caption}
\usepackage{fancyhdr}
\usepackage{emptypage}
\usepackage{indentfirst}
\usepackage{ragged2e}
\usepackage{booktabs}
\usepackage{array}
\usepackage{amsmath}
\usepackage{hyperref}
\usepackage{float}
\usepackage{longtable}

%--------------------------------------------------------------
% TIKZ + DIAGRAM PACKAGES
%--------------------------------------------------------------
\usepackage{tikz}
\usepackage{pgfgantt}
\usetikzlibrary{shapes.geometric,arrows.meta,positioning,fit,
                backgrounds,calc,decorations.pathreplacing,matrix}
\usepackage{xcolor}

% Shared TikZ styles reused across diagrams
\tikzset{
  box/.code       ={\pgfkeysalso{draw=#1,fill=#1,fill opacity=0.12,rounded corners=5pt,
                   align=center,font=\small}},
  modbox/.code    ={\pgfkeysalso{draw=#1,fill=#1,fill opacity=0.12,rounded corners=4pt,
                   align=center,font=\scriptsize,inner sep=4pt}},
  entity/.style  ={draw,thick,fill=blue!10,rounded corners=4pt,
                   align=center,font=\small,minimum width=2.0cm,
                   minimum height=0.8cm},
  process/.style ={draw,thick,fill=green!10,circle,
                   align=center,font=\scriptsize,
                   minimum size=1.4cm},
  dstore/.style  ={draw,fill=yellow!15,minimum width=2.4cm,
                   minimum height=0.7cm,align=center,font=\scriptsize},
  arr/.style     ={->,>=Stealth,thick},
  darr/.style    ={<->,>=Stealth,thick,gray!60},
}

%--------------------------------------------------------------
% BLANK LINE HELPER
%--------------------------------------------------------------
\newcommand{\blanklines}[1]{\vspace{#1\dimexpr 18pt\relax}}

%--------------------------------------------------------------
% PAGE GEOMETRY
%--------------------------------------------------------------
\geometry{
  a4paper,
  top=1.2in,
  bottom=1.0in,
  left=1.2in,
  right=1.0in,
  headheight=14pt,
  headsep=0.2in,
  footskip=0.4in
}

%--------------------------------------------------------------
% LINE SPACING
%--------------------------------------------------------------
\setstretch{1.5}

%--------------------------------------------------------------
% PARAGRAPH SPACING
%--------------------------------------------------------------
\setlength{\parskip}{6pt}
\setlength{\parindent}{0pt}

%--------------------------------------------------------------
% HEADER / FOOTER
%--------------------------------------------------------------
\pagestyle{fancy}
\fancyhf{}
\fancyhead[C]{\fontsize{12}{14}\selectfont}
\fancyfoot[C]{\thepage}
\renewcommand{\headrulewidth}{0pt}
\renewcommand{\footrulewidth}{0pt}

\fancypagestyle{plain}{
  \fancyhf{}
  \fancyfoot[C]{\thepage}
  \renewcommand{\headrulewidth}{0pt}
  \renewcommand{\footrulewidth}{0pt}
}

%--------------------------------------------------------------
% CHAPTER STYLE
% Use Roman numerals (CHAPTER I, II, III ...) per template
%--------------------------------------------------------------
\renewcommand{\thechapter}{\Roman{chapter}}
\titleformat{\chapter}[block]
  {\fontsize{16}{24}\bfseries\centering}
  {CHAPTER \thechapter}
  {1em}
  {}
\titlespacing*{\chapter}{0pt}{0pt}{24pt}

%--------------------------------------------------------------
% SECTION STYLE
%--------------------------------------------------------------
\titleformat{\section}
  {\fontsize{14}{21}\bfseries\raggedright}
  {\thesection}
  {1em}
  {}
\titlespacing*{\section}{0pt}{12pt}{6pt}

%--------------------------------------------------------------
% SUBSECTION STYLE
%--------------------------------------------------------------
\titleformat{\subsection}
  {\fontsize{12}{18}\bfseries\raggedright}
  {\thesubsection}
  {1em}
  {}
\titlespacing*{\subsection}{0pt}{12pt}{6pt}

%--------------------------------------------------------------
% TABLE OF CONTENTS
%--------------------------------------------------------------
\renewcommand{\contentsname}{Contents}
\renewcommand{\cfttoctitlefont}{\fontsize{16}{24}\bfseries\hfill}
\renewcommand{\cftaftertoctitle}{\hfill}
\setlength{\cftbeforetoctitleskip}{0pt}
\setlength{\cftaftertoctitleskip}{36pt}

\renewcommand{\cftchapfont}{\bfseries\fontsize{12}{18}\selectfont}
\renewcommand{\cftchappagefont}{\bfseries\fontsize{12}{18}\selectfont}
\setlength{\cftbeforechapskip}{6pt}

% Add "CHAPTER" prefix in TOC so entries read
% "CHAPTER I    Introduction    1" instead of "1 Introduction"
\renewcommand{\cftchappresnum}{CHAPTER\ }
\renewcommand{\cftchapaftersnum}{\quad}
\settowidth{\cftchapnumwidth}{CHAPTER\ XII\quad}

\renewcommand{\cftsecfont}{\fontsize{12}{18}\selectfont}
\renewcommand{\cftsecpagefont}{\fontsize{12}{18}\selectfont}
\setlength{\cftsecindent}{0.5in}
\setlength{\cftsubsecindent}{1.0in}

%--------------------------------------------------------------
% LIST OF TABLES
%--------------------------------------------------------------
\renewcommand{\listtablename}{List of Tables}
\renewcommand{\cftlottitlefont}{\fontsize{16}{24}\bfseries\hfill}
\renewcommand{\cftafterlottitle}{\hfill}
\setlength{\cftbeforelottitleskip}{0pt}
\setlength{\cftafterlottitleskip}{36pt}
\renewcommand{\cfttabfont}{\fontsize{12}{18}\selectfont}
\renewcommand{\cfttabpagefont}{\fontsize{12}{18}\selectfont}

%--------------------------------------------------------------
% LIST OF FIGURES
%--------------------------------------------------------------
\renewcommand{\listfigurename}{List of Figures}
\renewcommand{\cftloftitlefont}{\fontsize{16}{24}\bfseries\hfill}
\renewcommand{\cftafterloftitle}{\hfill}
\setlength{\cftbeforeloftitleskip}{0pt}
\setlength{\cftafterloftitleskip}{36pt}
\renewcommand{\cftfigfont}{\fontsize{12}{18}\selectfont}
\renewcommand{\cftfigpagefont}{\fontsize{12}{18}\selectfont}

%--------------------------------------------------------------
% FIGURE & TABLE CAPTION
%--------------------------------------------------------------
\captionsetup{
  font={rm,normalfont},
  labelfont=normalfont,
  textfont=normalfont,
  justification=centering,
  size=small,
  skip=6pt
}
\captionsetup[table]{position=above, skip=6pt}
\captionsetup[figure]{position=below, skip=6pt}

\renewcommand{\thefigure}{\arabic{chapter}.\arabic{figure}}
\renewcommand{\thetable}{\arabic{chapter}.\arabic{table}}

%--------------------------------------------------------------
% HYPERREF
%--------------------------------------------------------------
\hypersetup{
  colorlinks=false,
  pdfborder={0 0 0}
}

%==============================================================
% DOCUMENT BEGINS
%==============================================================
\begin{document}

%--------------------------------------------------------------
% COVER PAGE
%--------------------------------------------------------------
\thispagestyle{empty}
\begin{center}

  {\fontsize{14}{21}\selectfont CSE 3200: System Development Project}

  {\fontsize{18}{27}\bfseries\scshape
    Invigilo: Design and Development of a \\
    Secure Desktop Examination Platform  \\
  }

  \blanklines{2}

  {\fontsize{14}{21}\selectfont By}

  \blanklines{3}

  {\fontsize{14}{21}\bfseries Md. Nafiz Ahmed}\\[0pt]
  {\fontsize{14}{21}\selectfont Roll: 2107001}

  {\fontsize{14}{21}\bfseries \&}

  {\fontsize{14}{21}\bfseries Dewan Salman Rahman Zisan}\\[0pt]
  {\fontsize{14}{21}\selectfont Roll: 2107015}

  \vfill

% Logo
  \IfFileExists{logo.png}{%
    \includegraphics[width=1.14in,height=1.0in]{logo.png}%
  }{%
    \fbox{\parbox[c][1.0in][c]{1.14in}{\centering KUET\\Logo}}%
  }

  \blanklines{2}

  \vfill
  \begin{spacing}{1.2}
    {\fontsize{12}{14}\bfseries
      Department of Computer Science and Engineering\\
      Khulna University of Engineering \& Technology\\
      Khulna 9203, Bangladesh\\
      April, 2026
    }
  \end{spacing}

\end{center}

%--------------------------------------------------------------
% TITLE PAGE (page i)
%--------------------------------------------------------------
\newpage
\pagenumbering{roman}
\setcounter{page}{1}
\thispagestyle{plain}
\addcontentsline{toc}{chapter}{Title Page}

\begin{center}

  {\fontsize{18}{27}\bfseries
     Invigilo: Design and Development of a \\
    Secure Desktop Examination Platform  \\
  }

  \blanklines{3}

  {\fontsize{12}{18}\selectfont By}

  \blanklines{2}

  {\fontsize{12}{18}\bfseries Md. Nafiz Ahmed}\\[0pt]
  {\fontsize{12}{18}\selectfont Roll: 2107001}

  {\fontsize{12}{18}\bfseries \&}

  {\fontsize{12}{18}\bfseries Dewan Salman Rahman Zisan}\\[0pt]
  {\fontsize{12}{18}\selectfont Roll: 2107015}

\end{center}

\blanklines{4}

\noindent{\fontsize{12}{18}\bfseries Supervisor:}\\[0pt]
\hspace*{0.8in}{\fontsize{12}{18}\bfseries Waliul Islam Sumon}\\[0pt]
\hspace*{0.8in}{\fontsize{12}{18}\selectfont Lecturer}\\[0pt]
\hspace*{0.8in}{\fontsize{12}{18}\selectfont Department of Computer Science and Engineering}\\[0pt]
\hspace*{0.8in}{\fontsize{12}{18}\selectfont Khulna University of Engineering \& Technology}
\hfill \underline{\hspace{1.5in}}\\
\hspace*{0.8in}
\hfill {\fontsize{12}{18}\selectfont Signature}

\vfill

\begin{center}
  {\fontsize{12}{18}\selectfont
    Department of Computer Science and Engineering\\
    Khulna University of Engineering \& Technology\\
    Khulna 9203, Bangladesh\\
    April, 2026
  }
\end{center}

%--------------------------------------------------------------
% ACKNOWLEDGMENT (page ii)
%--------------------------------------------------------------
\newpage
\setcounter{page}{2}
\thispagestyle{plain}
\addcontentsline{toc}{chapter}{Acknowledgment}

\begin{center}
  {\fontsize{16}{24}\bfseries Acknowledgment}
\end{center}

\blanklines{2}

\begin{justify}
  {\fontsize{12}{18}\selectfont
    To begin with, we are grateful to Allah Almighty for blessing us with patience, a sound mind and strength which enabled us to successfully take this project from its inception stage to the development of a fully functional desktop application within the limited timeframe of six months.

     First and foremost, we would like to thank our honorable supervisor who provided us with valuable inputs during our weekly meetings which helped us shape up the Invigilo project to what we see now. The questions that our supervisor asked us regarding the invigilation of the exams and how teachers would be able to supervise the students through this technology forced us to rethink the approach we had taken to develop the proctoring layer.

    It would be good to thank the faculty members of the Department of Computer Science \& Engineering of KUET who introduced us to the topics related to the Operating System, Database Management System, Software Engineering, and Computer Networking which provided us with an understanding of many important concepts used to develop Invigilo, such as, JWT based session handling, designing relational schema, and process isolation.

   Our friends and class mates are thankful to us for acting as beta testers in our lab where their feedback on the handling of room codes, waiting room interface, and coding question editor interface was really invaluable to us in finding out problems relating to usability of the application.
  }
\end{justify}

\vspace{12pt}
\begin{flushright}
  {\fontsize{12}{18}\bfseries Authors}
\end{flushright}

%--------------------------------------------------------------
% ABSTRACT (page iii)
%--------------------------------------------------------------
\newpage
\setcounter{page}{3}
\thispagestyle{plain}
\addcontentsline{toc}{chapter}{Abstract}

\begin{center}
  {\fontsize{16}{24}\bfseries Abstract}
\end{center}

\blanklines{2}

\begin{justify}
  {\fontsize{12}{18}\selectfont
    Assessment using online means at Bangladeshi universities continues to rely extensively on browser-based quiz modules and form-type applications that lack effective ways of ensuring test validity. It is easy for a student to toggle to a new tab, start a chat app, launch another browser window, or even allow the other person to view what is on his/her screen without the quiz module being able to pick up any anomaly. On the contrary, teachers lack a common dashboard showing them which students have logged in, who faces problems with the internet connection, who has kept changing focus, or who should be force-submitted immediately.

    In this document, we describe Invigilo, which is a secure desktop examination software that will be used for lab quizzes and class tests in universities. Invigilo has been developed using Electron 39 for the desktop container, React 19 with Vite for the front end UI, Node.js with Express for back end API, and PostgreSQL for persistence. Role based dashboards implemented include Admin dashboard, which will be used for creating accounts as well as bulk uploads of students' details, Teacher dashboard, which allows teachers to create and monitor examinations, and the student dashboard, which provides students the ability to join an examination either via normal login or room code of six characters. Authentication process involves use of JSON Web Tokens and role based middleware, while password storage uses bcrypt.

    The desktop client locks itself into fullscreen always on top mode the moment an exam starts. The Electron main process actively blocks Alt+F4, F11 and most window manipulation shortcuts, monitors the focus state of the window, scans the running process list for forbidden tools such as remote desktop and screen sharing applications, and reports every violation back to the backend with a severity tag. A second monitoring layer runs inside the renderer using BlazeFace for browser based face detection. It captures snapshots every three seconds and emits events when no face, multiple faces, looking away or looking down conditions are detected. Teachers can view these events live, freeze any participant who is misbehaving, or force submit an active attempt with a single click.

    In an attempt to cater for classroom environments where students are frequently given coding homeworks, Invigilo offers several types of questions in one quiz/test submission, namely multiple choice, writing and programming question(s). The latter runs on the server side and will be evaluated in any one of the three possible languages: Javascript, Python or C++. There are restrictions on compile time and runtime of eight and twelve seconds respectively. In addition, question order will always be randomized but can be repeated in case of disputes.

    Quantitatively, the results for the current version show that there are 39 API endpoints in the backend, distributed across seven modules, seven core database tables, and a renderer component that runs in 11 to 13 seconds, depending on the host computer. There are also 12 test cases manually written for the integration of all aspects of the exam process in the current version. The limitations imposed by the current version of the software include no support for containerization, dependence on a host compiler, and the inability to perform automated performance testing for very large class sizes.
  }
\end{justify}

%--------------------------------------------------------------
% TABLE OF CONTENTS
%--------------------------------------------------------------
\newpage
\setcounter{page}{4}
\thispagestyle{plain}
\addcontentsline{toc}{chapter}{Contents}
\tableofcontents

%--------------------------------------------------------------
% LIST OF TABLES
%--------------------------------------------------------------
\newpage
\thispagestyle{plain}
\addcontentsline{toc}{chapter}{List of Tables}
\listoftables

%--------------------------------------------------------------
% LIST OF FIGURES
%--------------------------------------------------------------
\newpage
\thispagestyle{plain}
\addcontentsline{toc}{chapter}{List of Figures}
\listoffigures

%--------------------------------------------------------------
% MAIN MATTER
%--------------------------------------------------------------
\newpage
\pagenumbering{arabic}
\setcounter{page}{1}

%==============================================================
% CHAPTER I: Introduction
%==============================================================
\chapter{Introduction}

\section{Introduction}

\begin{justify}
  {\fontsize{12}{18}\selectfont
    Examinations are the core feedback mechanism in a university. Their results dictate the students' grade, scholarship, admission chances, and very often the type of job he can get after graduating. With so much at stake, honesty becomes just as important as the knowledge that is tested. In the traditional pen-and-paper format of an examination, this is achieved through invigilation, seating plan, printed questions, and secure answer scripts. The cost of cheating here is relatively low – the institution has to give out a new exam and face the embarrassment of having its schedule disrupted.

    Several Bangladeshi universities like KUET have resorted to using online examinations during class tests, laboratory tests, and even mid-term exams. There have been attempts to use online forms and quizzes provided by Google, Microsoft, Moodle quiz modules, as well as self-hosted systems based on PHP and a question bank. They successfully achieved what needed to be done – questions were asked, answers were submitted. However, all of these platforms were not designed with exam invigilation in mind. A website tab is inherently friendly to the user. It does not prevent him from changing tabs, opening new windows, contacting his friends via messenger, or running an IDE side-by-side. All that the teacher can do is make assumptions based on timestamps, missed answers, and his own intuition.

    This particular project, called Invigilo, was conceived with this very gap in mind. We do not treat the test paper as a regular web page, but rather as a controlled environment for executing the test. The candidate does not fire up his/her browser to access the test website; instead, the candidate fires up a native application which, from the very moment the test begins, hijacks the entire screen, disables all possible exit points, monitors the process list, and monitors the webcam feed. The instructor does not sit before an ever-repeating list of names; he/she sits before a live participant panel, complete with status changes, violations, and snapshots.

    In addition to control and monitoring, the second reason behind the design of Invigilo was a matter of practicality due to a CSE lab quiz. A normal lab quiz comprises questions ranging from multiple choice ones for theories to explanation and coding questions. In order to conduct quizzes using existing online proctoring solutions, teachers will either need to use three separate systems or make do without certain types of questions. The system Invigilo aims to provide the flexibility to integrate all kinds of questions within one examination system.

   The third reason had to do with institutional control and security issues. A commercial proctoring solution raises suspicions when it comes to the location of camera recordings, the person in charge of such files, their storage times and the model behind detecting any suspicious activity. As Invigilo will operate in a local setting where the machine hosting the software belongs to the same institute as the invigilator, everything is kept locally: snapshots are bound in size, event records are structured in rows and retention period configurable.
  }
\end{justify}

%--- Figure 1.1: Non-technical overview ---
\begin{figure}[t]
\centering
\begin{tikzpicture}[node distance=0pt]
  % ---- colour aliases ----
  \colorlet{admc}{blue!65!black}
  \colorlet{tchc}{green!55!black}
  \colorlet{stuc}{orange!75!black}
  \colorlet{sysc}{violet!65!black}

  % ---- central system box ----
  \node[draw=sysc,fill=sysc!8,rounded corners=10pt,
        minimum width=5.4cm,minimum height=3.8cm,
        label={[font=\bfseries\small,text=sysc]center:}] (sys) at (5.5,0) {};
  \node[font=\bfseries,text=sysc]        at (5.5, 1.15) {INVIGILO};
  \node[font=\scriptsize,text=sysc!80]   at (5.5, 0.65) {Secure Exam Platform};
  \node[draw=sysc!40,fill=sysc!5,rounded corners=3pt,
        font=\tiny,align=center,inner sep=3pt] at (4.4, 0.0)  {Desktop\\Lockdown};
  \node[draw=sysc!40,fill=sysc!5,rounded corners=3pt,
        font=\tiny,align=center,inner sep=3pt] at (6.6, 0.0)  {Webcam\\Proctoring};
  \node[draw=sysc!40,fill=sysc!5,rounded corners=3pt,
        font=\tiny,align=center,inner sep=3pt] at (5.5,-0.75) {Code Runner};

  % ---- role nodes ----
  \node[draw=admc,fill=admc!12,rounded corners=7pt,
        minimum width=2.8cm,minimum height=1.3cm,
        align=center,font=\small\bfseries,text=admc] (admin) at (5.5, 4.2)
        {ADMIN\\{\scriptsize User Management}};

  \node[draw=tchc,fill=tchc!12,rounded corners=7pt,
        minimum width=2.8cm,minimum height=1.3cm,
        align=center,font=\small\bfseries,text=tchc] (teacher) at (0.5, 0)
        {TEACHER\\{\scriptsize Exam Creation \& Monitoring}};

  \node[draw=stuc,fill=stuc!12,rounded corners=7pt,
        minimum width=2.8cm,minimum height=1.3cm,
        align=center,font=\small\bfseries,text=stuc] (student) at (5.5,-4.2)
        {STUDENT\\{\scriptsize Exam Attempt \& Results}};

  % ---- arrows ----
  \draw[arr,draw=admc,line width=1.2pt]
      (admin.south) -- (sys.north)
      node[midway,right,font=\scriptsize]{Create / Manage Users};

  \draw[arr,draw=tchc,line width=1.2pt]
      (teacher.east) -- (sys.west)
      node[midway,above,font=\scriptsize]{Exams, Questions, Controls};

  \draw[arr,draw=stuc,line width=1.2pt]
      (student.north) -- (sys.south)
      node[midway,right,font=\scriptsize]{Join, Answer, Submit};

  % ---- output labels from system ----
  \node[font=\scriptsize,text=admc!80,align=center]
      at (8.5, 2.6) {Stats \&\\Reports};
  \draw[->,>=Stealth,draw=admc!50,dashed]
      ([xshift=6pt]sys.north east) to[out=60,in=250] (8.5,2.0);

  \node[font=\scriptsize,text=tchc!80,align=center]
      at (8.5, 0.0) {Live\\Status};
  \draw[->,>=Stealth,draw=tchc!50,dashed]
      ([xshift=6pt]sys.east) -- (7.8,0.0);

  \node[font=\scriptsize,text=stuc!80,align=center]
      at (8.5,-2.6) {Evaluation\\Results};
  \draw[->,>=Stealth,draw=stuc!50,dashed]
      ([xshift=6pt]sys.south east) to[out=300,in=110] (8.5,-2.0);
\end{tikzpicture}
\caption{Non-technical overview of the Invigilo platform and its three user roles.}
\label{fig:1.1}
\end{figure}

\section{Problem Statement}

\begin{justify}
  {\fontsize{12}{18}\selectfont
    Despite the wide adoption of online assessment, the existing tooling used in most Bangladeshi universities falls short on several practical dimensions. The problems we identified before starting the project are listed below.

    \begin{itemize}
        \item \textbf{Browser Based Execution Without Lockdown:} Most quiz tools run inside a normal browser tab. There is nothing stopping the student from minimising the tab, opening a second browser, looking up answers on a phone or sharing the screen over Discord. Even if the tool detects a tab change, by the time the event is logged the damage is already done.

        \item \textbf{Missing Live Teacher Controls:} Once a quiz has started in a tool like Google Forms, the teacher has almost no way to intervene per student. The teacher cannot freeze a single participant, force submit only the suspicious ones, or push a warning to a particular roll number. Everything is all or nothing.

        \item \textbf{Fragmented Question Types:} Multiple choice, written answers and coding tasks are usually handled by three different platforms. The teacher has to combine the marks manually after the exam, and the student has to switch tools mid exam, which by itself becomes a focus violation in stricter setups.

        \item \textbf{No Built In Proctoring Signals:} Existing free tools do not monitor the camera, the running process list or the active window. A student may have a remote desktop session open in the background and the tool will never know.

        \item \textbf{Weak Identity Binding:} A roll number is usually verified once at login and then forgotten. There is no continuous evidence that the same person who logged in is the one actually answering the questions.

        \item \textbf{Inconsistent Code Execution Environment:} When a coding question is asked through an external editor, every student runs the code on their own machine with their own compiler version. Two students may submit the same answer and get different outputs because of environment differences.

        \item \textbf{Onboarding Friction for Real Classroom Use:} If a student forgets a password five minutes before the quiz starts, the entire exam can stall. Existing tools do not have a quick fallback like a room code that the teacher can share orally.

        \item \textbf{Audit Trail Gaps:} When a dispute happens after results are published, there is rarely a structured log that can be reviewed. The teacher can only point to a screenshot or a memory of what happened during the exam.

        \item \textbf{Privacy Concerns Around Commercial Proctoring:} Cloud based proctoring services collect video and behavioural data that the institution does not control. For a department running internal lab quizzes, this is a heavy and often unjustifiable trade off.
    \end{itemize}

    The core problem this project addresses can therefore be stated as follows. \textit{How do we design and implement a self hosted, role aware, desktop based examination platform that can deliver multiple choice, written and coding questions in a single integrated session, while giving teachers live participant level control and recording auditable proctoring evidence, without depending on external commercial services?}
  }
\end{justify}

\section{Objectives}

\begin{justify}
  {\fontsize{12}{18}\selectfont
    With the problem statement in mind, the specific objectives that we set for the project are the following.

    \begin{enumerate}
        \item Create three role based architecture (Admin, Teacher, Student) with all the API end points secured with role based middleware.
        \item Create a desktop client app for Electron that keeps the screen in fullscreen mode and on top during active exam attempts.
        \item Check for any violation, such as fullscreen mode exit, lack of focus, blocking keyboard shortcut, multi display usage and background process restrictions from the Electron main process.
        \item Create a single exam structure that can support multiple choice, written answer and coding question type.
        \item Execute the student submitted code in JavaScript, Python and C++ language on backend side with time restriction in execution and compiling.
        \item Provide a joining flow of six characters room code where a teacher can onboard students in the lab through verbal communication.
        \item Check whether a person's face is present on camera, number of faces present, face looking down or not with face detection using BlazeFace algorithm in the renderer.
        \item Give control to teacher for each individual participant including freeze, unfreeze and force submit.
        \item To provide a manual evaluation desk for written and coding answers that recalculates the total score automatically.
        \item To ensure that all submissions, violations, snapshots and evaluation records are persisted in a relational database for later audit.
        \item To document the entire system in a way that supports academic defense and future extension by another group.
    \end{enumerate}
  }
\end{justify}

\section{Scope}

\begin{justify}
  {\fontsize{12}{18}\selectfont
    The scope of the Invigilo project, in its current version, is described below in terms of what the system covers and what it deliberately leaves out.

    \begin{enumerate}
        \item The desktop client is built for Windows as the primary target. The same Electron codebase can be packaged for macOS and Linux, but the forbidden process scanner is currently tuned for Windows process names.
        \item The user interface is implemented in React 19 with Vite 7 as the build tool. Tailwind CSS is used for styling.
        \item The backend is implemented in Node.js with Express 4. Authentication uses jsonwebtoken with role based middleware, and password hashing uses bcryptjs.
        \item The persistent store is PostgreSQL. The schema is created at server boot with seven core tables and additive migration scripts for incremental columns such as the room code field.
        \item The face detection model is BlazeFace shipped through TensorFlow.js. Snapshots are sent to the backend at three second intervals when the student has consented to webcam usage.
        \item The coding execution engine supports Node.js for JavaScript, the system Python interpreter for Python, and g++ for C++. Each language has its own timeout configuration.
        \item The administrative scope includes single account creation, bulk student upload via CSV text, user activation and deactivation, and platform statistics.
        \item The system is not designed to replace high stakes board examinations. It is designed for university lab quizzes, class tests and internal departmental assessments.
        \item Out of scope in the current version are biometric identity verification, offline first local sync, fully containerised code sandboxing using Docker, and a mobile companion app.
    \end{enumerate}
  }
\end{justify}

\section{Unfamiliarity of the Problem/Topic/Solution}

\begin{justify}
  {\fontsize{12}{18}\selectfont
    Despite the familiarity of the topic of online quizzing itself, the combination of conditions that Invigilo solves is not readily available for purchase as a ready-made product for easy installation and deployment. Most commonly, the product will excel in solving some small part of the task, while ignoring others. A Learning Management System will allow managing the questions bank effectively but won't be able to lock the screen of a participant's computer. On the other hand, a lockdown browser will block the desktop, but it will lack functionality allowing the instructor to pause an individual student and submit them, proceeding with the test of the rest of the participants. An expensive commercial proctoring service will monitor the student via their webcam, but it will not allow specifying a question that needs to be solved by writing code in a specific language with a custom time limit.

    In terms of the local context of a department of engineering education in Bangladesh, the level of unfamiliarity will be more pronounced. Lab instructors whom I have informally polled have not used any solution that allows freezing the screen of a disruptive student during a test, forcefully submitting them and proceeding with the rest of the test participants without stopping. The feature of having a room code consisting of six characters that will enable a participant to enter a test without logging in using their credentials is also novel for most available products that the students are accustomed to.

    From the engineering perspective, the novel elements resulted from the composition of all parts rather than from any individual part. Individual technologies such as Electron, React, Express and PostgreSQL are each well documented systems on their own. The real challenge involved communicating between the Electron main process and the React renderer through an intermediate preload layer; the synchronization of the desktop locking process with the examination states in the backend; the handling of proctoring events without overwhelming the participant list monitoring process; and the securing of the execution endpoint while running on a bare VM without any kind of isolation container.

    The approach we have adopted for this assignment – namely, the use of a self hosted desktop application for proctoring with teacher controls – is certainly not the common one in a university environment laboratory. Awareness of that early in development allowed us to define realistic requirements and document design choices for potential reuse by another team.
  }
\end{justify}

\section{Project Planning and Work Distribution}

\begin{justify}
  {\fontsize{12}{18}\selectfont
    The project started in late December 2025 and ended in April 2026 for a duration of around seventeen weeks. We used the iterative agile model with iterations of about two to three weeks each. At the end of every iteration, we held an internal demo, created a list of open issues in a shared Google Document, and planned for the next iteration. The history of the Git repository, which shows contributions of thirty commits by each author out of two, bears testimony to the iterative nature of our development process.

%--- Figure 1.2: SDLC / Iterative Development Cycle ---
\begin{figure}[t]
\centering
\begin{tikzpicture}[
  phase/.code={\pgfkeysalso{draw=#1,fill=#1,fill opacity=0.15,rounded corners=6pt,
                minimum width=2.6cm,minimum height=1.0cm,
                align=center,font=\small\bfseries,text=black}},
  carrow/.code={\pgfkeysalso{->,>=Stealth,line width=1.4pt,draw=#1}}
]
  % Define 6 phases in a hexagonal ring
  % Angles: 90,30,330,270,210,150
  \def\R{3.2}
  \foreach \ang/\col/\lbl/\sub in {
      90/blue/{Requirements}{Role analysis, stack selection},
      30/green!60!black/{Design}{Architecture, schema, API},
      330/orange/{Implement}{Backend, frontend, security},
      270/red!70!black/{Test}{Integration runs, demo sessions},
      210/violet/{Review}{Feedback, issue logging},
      150/teal/{Deploy\,/\,Document}{Build check, report writing}
  }{
    \node[phase=\col,text width=2.4cm] (\ang)
        at (\ang:\R) {\lbl\\[2pt]{\tiny\color{black}\sub}};
  }

  % Circular arrows between phases
  \foreach \from/\to/\col in {90/30/blue,30/330/green!60!black,330/270/orange,270/210/red!70!black,210/150/violet,150/90/teal}{
    \draw[carrow=\col]
        (\from) to[bend left=28] (\to);
  }

  % Center label
  \node[font=\bfseries\small,align=center,text=gray!70!black] at (0,0)
      {Iterative\\Development\\Cycle};
\end{tikzpicture}
\caption{Iterative agile development cycle followed during the Invigilo project.}
\label{fig:1.2}
\end{figure}

    The tasks were divided mainly according to layers and expertise of the respective members. Md. Nafiz Ahmed was responsible for the desktop application skeleton, front-end dashboards, user interface adjustments arising out of actual use in the lab environment, and the documentation of the entire project. Dewan Salman Rahman Zisan was responsible for the back-end authentication mechanism, room code generation engine, security layer in the Electron main process, API for running coding tasks, AI assistant implementation, and webcam face detection.
  }
\end{justify}

\subsection{Work Distribution}

\begin{justify}
  {\fontsize{12}{18}\selectfont
    The contributions of the two team members, derived from the actual Git commit history of the repository, are listed below.

    \textbf{Developer 1: Md. Nafiz Ahmed (Roll: 2107001)}
    \begin{enumerate}
        \item Initial Electron application bootstrap and fullscreen window setup.
        \item Initial login screen and exam flow skeleton in the renderer.
        \item Modern Tailwind based login user interface without breaking the exam flow.
        \item Teacher dashboard with exam and question management interface.
        \item Replacement of all blocking native alert and confirm dialogs with custom in app modals to fix keyboard input focus issues in Electron fullscreen mode.
        \item Post submission flow improvements and prevention of re entry into already submitted exams (frontend side).
        \item Inline form refactor for question management instead of modal driven flow.
        \item Integration of the login user interface with the backend authentication API.
        \item User interface fixes for layout, repeated auto submit alerts, violation count display, violation severity warnings and environment warnings.
        \item Display of forbidden process violations to the student during an active exam.
        \item Coding interface user interface fixes from real usage feedback.
        \item Maintenance of internal team guides (NAFIZ\_GUIDE, DEWAN\_GUIDE) and documentation updates after every major milestone.
    \end{enumerate}

    \textbf{Developer 2: Dewan Salman Rahman Zisan (Roll: 2107015)}
    \begin{enumerate}
        \item Backend authentication system using JSON Web Tokens and PostgreSQL with bcrypt password hashing.
        \item Backend APIs for teacher exam management including create, update and lifecycle control endpoints.
        \item Database migration and backend implementation of the room code joining system.
        \item Frontend implementation of the room code system and the exam taking interface used by students.
        \item Electron main process security layer including fullscreen enforcement, blocking of Alt+F4 and F11, fullscreen exit detection, window focus loss detection, and z order enforcement to keep the exam window topmost.
        \item Detection of multiple displays and environment changes during an active exam.
        \item Forbidden background process detection with immediate termination logic and an instant scan path.
        \item Single auto submit enforcement and exam state reset after submission.
        \item Manual exam submit flow and exam only always on top behaviour.
        \item Webcam based face detection integration using BlazeFace and the snapshot upload pipeline.
        \item Coding interface with editor, language selection, run support and AI assistant flow updates.
        \item Backend logging and lifecycle stability fixes for the Electron main process.
    \end{enumerate}

    Although the layers are split, the two developers reviewed each other's pull requests before merging into the develop branch. Several commits, especially those tagged as integration fixes, are the result of pair debugging sessions where one author drove the keyboard and the other read the logs.
  }
\end{justify}

\subsection{Gantt Chart}

\begin{justify}
  {\fontsize{12}{18}\selectfont
    The project schedule was tracked at the granularity of weeks. Table 1.1 shows the indicative phase plan that we followed across the seventeen weeks of active development.
  }
\end{justify}

\begin{table}[H]
  \caption{Indicative Gantt Style Phase Plan for Invigilo}
  \label{tab:1.1}
  \centering
  {\fontsize{11}{14}\selectfont
    \begin{tabular}{|p{1.2in}|c|p{3in}|}
      \hline
      \textbf{Phase} & \textbf{Weeks} & \textbf{Core Deliverables} \\
      \hline
      Requirement and risk analysis & 1--2 & Problem framing, role identification, threat assumptions, stack selection \\
      \hline
      Architecture and schema design & 3--4 & Three tier architecture, ER diagram, route design, Electron preload bridge \\
      \hline
      Core backend implementation & 5--7 & Auth APIs, exam APIs, submission and evaluation APIs, room code engine \\
      \hline
      Frontend role dashboards & 8--10 & Student, Teacher and Admin dashboards in React \\
      \hline
      Security and proctoring integration & 11--12 & Desktop lock, violation pipeline, webcam face detection \\
      \hline
      Coding engine and AI assistant & 13--14 & Multi language code runner, teacher AI assistant flow \\
      \hline
      Integration and debugging & 15--16 & Cross module validation, focus and submission edge cases \\
      \hline
      Final verification and documentation & 17 & Build verification, manual test runs, report and slides \\
      \hline
    \end{tabular}
  }
\end{table}

%--- Figure 1.3: Gantt Chart ---
\begin{figure}[t]
\centering
\begin{tikzpicture}[x=0.72cm,y=0.52cm]
  % Grid and week headers
  \foreach \w in {1,...,17}{
    \draw[gray!25,thin] (\w,0) -- (\w,-8.5);
    \node[font=\tiny,gray!70] at (\w+0.5,0.35) {\w};
  }
  \draw[gray!25,thin] (18,0) -- (18,-8.5);
  \node[font=\scriptsize\bfseries,gray!60] at (9.5,0.9) {Week};

  % Phase definitions: {label}{start}{end}{color}{row}
  \foreach \lbl/\s/\e/\col/\row in {
    {Req.\ \& Risk Analysis}/1/2/blue!60/1,
    {Architecture \& Schema}/3/4/teal!70/2,
    {Core Backend}/5/7/green!55!black/3,
    {Frontend Dashboards}/8/10/orange!70/4,
    {Security \& Proctoring}/11/12/red!60/5,
    {Coding Engine \& AI}/13/14/violet!70/6,
    {Integration \& Debug}/15/16/brown!60/7,
    {Verification \& Docs}/17/17/gray!60/8
  }{
    % Row label
    \node[font=\scriptsize,anchor=east] at (0.9,-\row+0.5) {\lbl};
    % Bar
    \fill[fill=\col,fill opacity=0.4,draw=\col,rounded corners=2pt]
        (\s,-\row+0.12) rectangle (\e+1,-\row+0.88);
  }

  % Horizontal row separators
  \foreach \r in {0,...,8}{
    \draw[gray!20,thin] (1,-\r) -- (18,-\r);
  }
  % Outer border
  \draw[gray!50] (1,0) rectangle (18,-8);
\end{tikzpicture}
\caption{Gantt chart of the seventeen-week Invigilo development schedule.}
\label{fig:1.3}
\end{figure}

\subsection{RACI Matrix}

\begin{justify}
  {\fontsize{12}{18}\selectfont
    To make responsibility unambiguous we used a simple RACI framework where R stands for Responsible, A for Accountable, C for Consulted and I for Informed. Table 1.2 summarises the responsibilities for the major activities of the project.
  }
\end{justify}

\begin{table}[H]
  \caption{RACI Matrix for Invigilo Project Activities}
  \label{tab:1.2}
  \centering
  {\fontsize{11}{14}\selectfont
    \begin{tabular}{|p{2in}|c|c|c|}
      \hline
      \textbf{Activity} & \textbf{Nafiz} & \textbf{Zisan} & \textbf{Supervisor} \\
      \hline
      System architecture and schema design & A/R & C & I \\
      \hline
      Backend authentication and exam APIs & C & R & I \\
      \hline
      Electron main process security layer & C & A/R & I \\
      \hline
      Frontend dashboards (Student/Teacher/Admin) & A/R & C & I \\
      \hline
      Room code joining flow & C & R & I \\
      \hline
      Webcam proctoring integration & C & R & I \\
      \hline
      Coding execution engine & C & R & I \\
      \hline
      Manual integration testing & R & R & C \\
      \hline
      Report writing and slide preparation & R & C & C \\
      \hline
      Final demo and viva preparation & R & R & I \\
      \hline
    \end{tabular}
  }
\end{table}

\section{Applications of the Work}

\begin{justify}
  {\fontsize{12}{18}\selectfont
    The applications of Invigilo go beyond the immediate context of a single department. The system was designed for a CSE lab quiz, but the architecture is general enough to be reused in several scenarios.

    \begin{itemize}
        \item \textbf{Lab Quizzes for Programming Courses:} The most direct application. A typical CSE 2100 or CSE 3100 lab class has thirty students writing a short coding task in forty minutes. Invigilo can host this end to end including the code execution and the manual evaluation.

        \item \textbf{Class Tests with Mixed Question Types:} Theory courses often need a mixture of MCQ and short answer questions. Invigilo allows both inside one exam without the teacher needing to use two tools.

        \item \textbf{Internal Departmental Assessments:} Recruitment tests for teaching assistants, project group selection tests and similar internal assessments can be conducted under the same secure setup.

        \item \textbf{Audit Oriented Evaluations:} Whenever an evaluation may need to be reviewed later, for example for a grade dispute, the violation events, snapshots and submission records stored by Invigilo serve as structured evidence.

        \item \textbf{Hybrid Examinations:} A subset of the class can sit in the lab while another subset writes from home. Both groups go through the same flow because the desktop application enforces the same controls everywhere.

        \item \textbf{Faculty Recruitment Coding Tests:} Departments that interview lab demonstrators or research assistants can use Invigilo to pose a real coding question and run it under proctored conditions.

        \item \textbf{Training Workshops with Assessment:} Short training programs that include a graded test at the end can use Invigilo to ensure that participants sit the test fairly even when participants come from different institutions.

        \item \textbf{Contextual Evaluation in Resource Limited Settings:} Because the system is self hosted on department infrastructure, it can be deployed in environments where a paid commercial proctoring subscription is not financially feasible.
    \end{itemize}
  }
\end{justify}

\section{Organization of the Project}

\begin{justify}
  {\fontsize{12}{18}\selectfont
    This paper will be organized in the following manner. Chapter 2 gives an overview of the current state of affairs on the field of the online assessments and puts Invigilo into context with regard to learning management system quizzes, form-based tests, proctoring software as well as the available solutions to the problems of cheating during examinations. Chapter 3 explains in detail the design approach that was used when creating the platform, its tri-level architecture, the employed data models, algorithms used for the generation of the room codes and randomization of the questions, and the processes followed by each type of role. Chapter 4 provides information about the implementation of the platform, including the experiments performed, the collected results of both quantitative and qualitative nature, manual integration testing cases and mapping the original set of goals against the accomplished achievements. Chapter 5 focuses on the socio-economic, ethic, legal and ecological factors related to the operation of a proctored examination platform. Chapter 6 talks about some interesting and challenging engineering problems that emerged while working on the project. Chapter 7 summarizes the work, listing the shortcomings of the existing version and suggesting some directions for improvement in the future.
  }
\end{justify}

%==============================================================
% CHAPTER II: Related Works
%==============================================================
\chapter{Related Works}

\section{Introduction}

\begin{justify}
  {\fontsize{12}{18}\selectfont
    Online testing is a well-established area. Software systems capable of presenting questions via the Internet and recording the answers have been around for over twenty years, dating back to rudimentary quizzes on Moodle to modern day cloud-based services with adaptive test banks. Prior to adopting the Invigilo architecture, we conducted an analysis of the current market offerings to establish the state of affairs in terms of solved versus unsolved problems in the context of an engineering course at an educational institution in Bangladesh.

    These systems can be classified into four distinct categories. The first category is the LMS quiz plugin, exemplified by Moodle Quiz and Canvas Quiz. The second category is the form-based lightweight software, with Google Forms and Microsoft Forms representing the mainstream options in this space. The third category is the lockdown browser together with LMS, with Respondus LockDown Browser and Safe Exam Browser being the most familiar products in this class. The fourth category consists of proprietary AI proctoring solutions, represented by ProctorU, ExamSoft, Proctorio and Honorlock among others.

    All of these approaches have their advantages and drawbacks. The question bank provided by learning management systems in the form of quiz modules is extensive, although the exam itself is performed via the regular browser of the student, which means that focus loss will need to be monitored via the unreliable browser API. Form centric solutions are simple to implement yet not tailored to high stakes tests. While lockdown browsers provide better environment control, they may be difficult to integrate with other technologies, usually come at a hefty price and require an additional license. Lastly, commercially available AI proctoring systems are excellent at what they do, though there are concerns regarding their cost and data protection for a public educational institution administering a quiz within its own laboratory.

    Another group of technologies, which is quite similar to the previous one but more specific in nature, revolves around coding competition platforms including, but not limited to, HackerRank, Codility and LeetCode contests. These technologies excel at assessing coding skills; nevertheless, these are not designed as comprehensive exam platforms for mixing multiple choice and written questions with code assessment.

    The Bangladesh case is no exception to ad hoc arrangements as well. For example, we found courses where the teacher sent out a paper through e mail, the students downloaded it and wrote their responses in a Word document before sending it back through the same medium. In addition, there were even cases where students shared their answers on WhatsApp during an ongoing quiz session. These solutions are not engineering competitors of Invigilo, but they do define the baseline reality which this project aims to improve upon.

    Another important consideration regarding the academic literature is that, since 2020, there has been an explosion in scholarly articles addressing secure online examinations, due largely to the necessity for remote learning brought about by the pandemic. In particular, various authors have studied the tradeoff between proctoring severity and student wellbeing, the rate of false positive detections when using automated cheater detection systems, and whether recorded examination footage could be considered legally admissible. The latter research was especially useful in formulating our policy-based sections, especially Chapter 5.
  }
\end{justify}

\section{Discussion}

\begin{justify}
  {\fontsize{12}{18}\selectfont
    Our review of related works confirmed that no single tool already provides the integrated mix of features that Invigilo offers. The closest analogue would be a paid combination of a learning management system, a lockdown browser, a commercial proctoring service and a coding judge platform, glued together by the institution's IT team. For a small department this is neither financially nor operationally realistic. Table 2.1 summarises the comparison across the categories that we surveyed.
  }
\end{justify}

\begin{table}[H]
  \caption{Comparative View of Online Assessment Approaches}
  \label{tab:2.1}
  \centering
  {\fontsize{10}{13}\selectfont
    \begin{tabular}{|p{1.6in}|c|c|c|c|c|}
      \hline
      \textbf{Criterion} & \textbf{LMS Quiz} & \textbf{Form Tool} & \textbf{Lockdown + LMS} & \textbf{AI Proctoring} & \textbf{Invigilo} \\
      \hline
      Role aware admin/teacher/student split & Partial & Limited & Partial & Limited & Full \\
      \hline
      Room code onboarding & Rare & Common & Rare & Rare & Full \\
      \hline
      Per student freeze and force submit & No & No & Limited & No & Yes \\
      \hline
      Built in coding execution & Plugin & No & Plugin & No & Yes \\
      \hline
      Desktop level fullscreen lock & No & No & Yes & Partial & Yes \\
      \hline
      Forbidden process scanning & No & No & Limited & Limited & Yes \\
      \hline
      Webcam face event reporting & No & No & Limited & Yes & Yes \\
      \hline
      Manual evaluation of subjective answers & Yes & Limited & Yes & Limited & Yes \\
      \hline
      Self hosted by institution & Yes & No & Partial & No & Yes \\
      \hline
      License cost & Free/Paid & Free & Paid & Paid & Free \\
      \hline
    \end{tabular}
  }
\end{table}

\begin{justify}
  {\fontsize{12}{18}\selectfont
    The table shows Invigilo as winning on each of these criteria, although we must not pretend that what we mean when we put a tick in this box is what a commercial vendor means. The implementation of the system includes all of the above to the extent that is sufficient for use in a university environment in lab quizzes and class tests; it does not provide all of these elements at the level of polish seen in mature software used in millions of exams worldwide and trained over many years. A commercial company providing AI-based proctoring services has many years worth of face detection training data behind their detection model, as well as many lawyers behind its data handling policies. Invigilo uses an open source face detection algorithm.

    Two additional lessons learned while reviewing the literature in the field of proctoring influenced our final decision of how to design Invigilo. Firstly, most papers mention students' right to know what information is being recorded and for what purpose. Invigilo makes it optional for students to use their webcams and this choice is available in the dashboard before any recording takes place. Secondly, the vast majority of successful lockdown solutions include process-level monitoring of potentially malicious software, and we adopted this strategy as well while keeping a small list of such processes transparent for students.

    Taken together, these design choices place Invigilo in a niche that is not well served by any single existing tool, namely a self hosted, transparent, mixed format examination platform with practical proctoring designed for the workflow of a Bangladeshi engineering department.
  }
\end{justify}

%==============================================================
% CHAPTER III: Methodology
%==============================================================
\chapter{Methodology}

\section{Introduction}

\begin{justify}
  {\fontsize{12}{18}\selectfont
    In this chapter, Invigilo is explained step by step from an architectural perspective. The chapter describes the system architecture, data model, key algorithms, the workflow for each actor, and the reasons for choosing this particular technology stack. Whenever possible, we use real configurations used in our implementation instead of some abstract representations because otherwise, we will have no choice but to describe what a reader would see if he/she were to clone our repository and run the application on his/her computer.

   The approach we took in designing Invigilo is guided by the principle of least surprise. We implemented each component in the manner which required the least number of assumptions about other components and maintained the narrowest interfaces between layers. The Electron main process communicates with the renderer using a single bridge called "preload". The renderer communicates with the backend via a single REST API client. The backend communicates with PostgreSQL through a single connection pool. It helped us easily debug and solve integration problems and allowed us to make changes in one layer without affecting others.
  }
\end{justify}

\section{Detailed Methodology}

\begin{justify}
  {\fontsize{12}{18}\selectfont
    Invigilo is an implementation of a three-tiered desktop application. The client tier comprises an Electron 39 application that bundles a Chromium renderer with React 19. The service tier comprises an Express-based Node.js server 4 that exposes a REST API secured by JWT middleware. The database is comprised of a PostgreSQL database that has seven main tables and some additive migration scripts for incremental column addition.

    Building upon the foundation provided by the above, we implemented two cross-cutting concerns as separate modules within our application. The first concern, a security module for Electron’s main process that handles examination mode, violation detection, and the forbidden process detector is implemented independent of the renderer. The proctoring concern, on the other hand, involves the use of BlazeFace from within the renderer to generate events that are then emitted to the backend via a proctoring route.

    Our software development methodology was an iterative one. Our cycles lasted two to three weeks, and at the end of each, an internal demo was performed on two laptops with one configured as the instructor and the other as the examinee. Bug fixes that emerged during these demos were tracked as small commits with descriptive messages, which is visible in the Git history under tags like \texttt{fix(security)}, \texttt{fix(frontend)}, \texttt{fix(backend)} and \texttt{feat(security)}.

    The actors of the system and their primary responsibilities are summarised below.

    \begin{itemize}
        \item \textbf{Admin:} Provisions teacher and student accounts either one at a time or through a CSV upload. Activates and deactivates users. Reviews high level statistics such as the number of teachers, students and exams.
        \item \textbf{Teacher:} Creates exams with metadata such as duration, type, presentation mode, and webcam requirement. Authors questions of three kinds: MCQ, written and coding. Shares the room code with students. Starts the exam. Monitors live participant status, violations and proctoring snapshots. Performs freeze, unfreeze and force submit actions per participant. Evaluates written and coding answers on the evaluation desk.
        \item \textbf{Student:} Logs in using either email or seven digit roll number with password, or joins through a six character room code. Waits in the waiting room until the teacher starts the exam. Answers questions in the configured presentation mode. Runs and tests code for coding questions. Submits the exam manually or has it submitted automatically on timeout or by force. Views detailed results after evaluation.
    \end{itemize}
  }
\end{justify}

\subsection{Problem Design and Analysis}

\begin{justify}
  {\fontsize{12}{18}\selectfont
    The problem of building a secure desktop examination platform was decomposed into smaller engineering subproblems that could be tackled independently and then integrated. Table 3.1 lists the major subproblems and the design choices made for each.
  }
\end{justify}

\begin{table}[H]
  \caption{Design Decomposition of the Invigilo System}
  \label{tab:3.1}
  \centering
  {\fontsize{10}{14}\selectfont
    \begin{tabular}{|p{1.5in}|p{1.7in}|p{2in}|}
      \hline
      \textbf{Subproblem} & \textbf{Design Decision} & \textbf{Rationale} \\
      \hline
      Desktop lockdown & Electron main process with fullscreen, always on top, blocked shortcuts & A real browser tab cannot block its own minimisation \\
      \hline
      Role separation & Three roles enforced both at route level and at resource ownership level & A single missed check should not let a student act as a teacher \\
      \hline
      Authentication & JSON Web Tokens with bcrypt hashed passwords & Stateless tokens make API scaling and testing easier \\
      \hline
      Room code joining & Six character upper case alphanumeric code with an ambiguity reduced alphabet & Easy to read aloud in a noisy lab and unique enough for the cohort size \\
      \hline
      Question randomisation & Hash seeded deterministic shuffle per student per exam & Fair to students and reproducible if a dispute arises \\
      \hline
      Mixed question types & Single submission record with type tagged answers & Lets the evaluation desk treat all question types uniformly \\
      \hline
      Coding execution & Spawned host processes with strict timeouts & Faster to implement than full container isolation while still bounded \\
      \hline
      Proctoring events & Throttled BlazeFace events plus periodic snapshots & Avoids flooding the backend with redundant events \\
      \hline
      Live teacher control & Polling based participant panel with intervention endpoints & Predictable and easier to debug than websocket flows for the project scope \\
      \hline
    \end{tabular}
  }
\end{table}

\subsection{Overall Framework and Flowchart}

\begin{justify}
  {\fontsize{12}{18}\selectfont
    The architecture diagram of Invigilo at a high level is presented in Figure 3.1 in the form of a textual block. The Electron application runs on the student’s computer. In the Electron container, the main process manages the controls at an OS level while the renderer process handles the React-based UI components. The communication between the React-based renderer process and the Express-based back-end takes place through HTTPS protocol.
  }
\end{justify}

%--- Figure 3.1: System Architecture ---
\begin{figure}[t]
\centering
\begin{tikzpicture}[x=1cm,y=1cm]
  % ---- Layer 1: Electron Client ----
  \node[draw=blue!60,fill=blue!6,rounded corners=7pt,
        minimum width=12.5cm,minimum height=2.1cm,
        label={[font=\small\bfseries,text=blue!70!black,anchor=west]
               north west:\strut Electron Desktop Client}]
       (client) at (6.25,5.2) {};
  \foreach \lbl/\cx in {
      {Main Process\\[-2pt]{\tiny lockdown, violations}}/2.0,
      {React Renderer\\[-2pt]{\tiny UI dashboards}}/4.8,
      {BlazeFace\\[-2pt]{\tiny face detection}}/7.6,
      {Monaco\\[-2pt]{\tiny code editor}}/10.2
  }{
    \node[modbox=blue,minimum width=2.2cm,minimum height=1.3cm] at (\cx,5.2) {\lbl};
  }

  % ---- connector 1 ----
  \draw[darr,blue!50,line width=1.2pt]
      (6.25,4.1) -- (6.25,3.55)
      node[midway,right,font=\scriptsize]{REST / JSON \quad IPC (preload)};

  % ---- Layer 2: Express Backend ----
  \node[draw=green!55!black,fill=green!5,rounded corners=7pt,
        minimum width=12.5cm,minimum height=2.1cm,
        label={[font=\small\bfseries,text=green!55!black,anchor=west]
               north west:\strut Express Backend (Node.js)}]
       (backend) at (6.25,2.6) {};
  \foreach \lbl/\cx in {
      {Auth}/1.3,
      {Exam}/2.9,
      {Question}/4.65,
      {Proctor}/6.45,
      {Admin}/8.2,
      {AI}/9.7,
      {Code}/11.2
  }{
    \node[modbox=green!55!black,minimum width=1.5cm,minimum height=1.2cm] at (\cx,2.6) {\lbl};
  }

  % ---- connector 2 ----
  \draw[darr,green!50!black,line width=1.2pt]
      (6.25,1.5) -- (6.25,0.95)
      node[midway,right,font=\scriptsize]{pg driver / SQL};

  % ---- Layer 3: Database ----
  \node[draw=orange!70,fill=orange!6,rounded corners=7pt,
        minimum width=12.5cm,minimum height=1.9cm,
        label={[font=\small\bfseries,text=orange!70!black,anchor=west]
               north west:\strut PostgreSQL Database}]
       (db) at (6.25,0.05) {};
  \foreach \lbl/\cx in {
      {Users}/1.6,
      {Exams}/3.4,
      {Questions}/5.3,
      {Submissions}/7.4,
      {Participants}/9.5,
      {Proct.\ Data}/11.3
  }{
    \node[modbox=orange,minimum width=1.8cm,minimum height=1.1cm] at (\cx,0.05) {\lbl};
  }

  % ---- External: Groq + Runtimes ----
  \node[draw=violet!60,fill=violet!8,rounded corners=5pt,
        minimum width=1.8cm,minimum height=1.4cm,align=center,font=\scriptsize]
       (groq) at (13.8,2.6) {Groq\\LLM API};
  \draw[arr,violet!60] (12.5,2.9) -- (groq.west)
      node[midway,above,font=\tiny]{HTTPS};

  \node[draw=brown!60,fill=brown!8,rounded corners=5pt,
        minimum width=1.8cm,minimum height=1.4cm,align=center,font=\scriptsize]
       (run) at (13.8,1.2) {Node/Python\\g++ Runtimes};
  \draw[arr,brown!60] (12.5,2.3) to[out=0,in=180] (run.west)
      node[midway,above,font=\tiny]{spawn};
\end{tikzpicture}
\caption{High-level three-tier architecture of the Invigilo system.}
\label{fig:3.1}
\end{figure}

\begin{justify}
  {\fontsize{12}{18}\selectfont
    The student attempt workflow is the most security sensitive flow in the system. Figure 3.2 shows the sequence of states a student goes through, from logging in to viewing the final result.
  }
\end{justify}

%--- Figure 3.2: Student Attempt Workflow ---
\begin{figure}[t]
\centering
\begin{tikzpicture}[
  step/.code={\pgfkeysalso{draw=#1,fill=#1,fill opacity=0.14,rounded corners=6pt,
               minimum width=5.6cm,minimum height=0.82cm,
               align=center,font=\small}},
  dec/.style={draw=orange!70,fill=orange!12,diamond,
              minimum width=3.2cm,minimum height=0.9cm,
              align=center,font=\scriptsize,aspect=2.8},
  fa/.style={arr,gray!60,line width=1.1pt},
  note/.style={font=\scriptsize,text=gray!70,anchor=west}
]
  \node[draw,fill=black!80,circle,text=white,font=\small,
        minimum size=0.5cm] (start) at (0,0) {};

  \node[step=blue]       (login)  at (0,-1.2)  {Login \textbf{or} Enter Exam (Room Key + Roll)};
  \node[step=blue!60]    (wait)   at (0,-2.4)  {Waiting Room --- polling every 3 s};
  \node[dec]             (started) at (0,-3.55) {Exam started?};
  \node[step=red!70]     (lock)   at (0,-4.8)  {Exam Mode Lock --- fullscreen, always on top};
  \node[step=green!55!black] (attempt) at (0,-6.0) {Answer Questions (MCQ / Written / Coding)};
  \node[dec]             (code)   at (0,-7.2)  {Coding question?};
  \node[step=teal!70]    (run)    at (3.5,-7.2) {Run Code, Check Output};
  \node[dec]             (sub)    at (0,-8.5)  {Submit condition?};
  \node[step=violet!70]  (done)   at (0,-9.7)  {Submission Persisted --- Lock Released};
  \node[step=gray!60]    (result) at (0,-10.9) {Result View with Per-Question Feedback};
  \node[draw,fill=black!80,circle,text=white,font=\small,
        minimum size=0.5cm] (end)  at (0,-12.1) {};

  % Vertical flow
  \foreach \from/\to in {start/login,login/wait,wait/started,
                          lock/attempt,attempt/code,done/result,result/end}{
    \draw[fa] (\from) -- (\to);
  }
  % Started decision
  \draw[fa] (started.south) -- (lock) node[midway,right,font=\scriptsize]{Yes};
  \draw[fa] (started.west)  to[out=180,in=180] (wait.west)
      node[near start,left,font=\scriptsize]{No};
  % Code decision
  \draw[fa] (code.east) -- (run) node[midway,above,font=\scriptsize]{Yes};
  \draw[fa] (run.south)  to[out=270,in=0] (sub.east);
  \draw[fa] (code.south) -- (sub) node[midway,right,font=\scriptsize]{No};
  % Submit decision
  \draw[fa] (sub.south) -- (done);
  \node[note] at (0.4,-9.1) {manual / timeout / force-submit};
\end{tikzpicture}
\caption{Student exam attempt workflow inside Invigilo.}
\label{fig:3.2}
\end{figure}

\begin{justify}
  {\fontsize{12}{18}\selectfont
    The major algorithms used inside this workflow are described below. The following algorithms are being used internally in the system and are coded in the backend controllers or Electron main process scripts respectively.

    \textbf{Room Code Generation.} A six letter alphanumeric code is generated based on a reduced ambiguous alphabet where letters zero, capital o, capital i and number one have been excluded. The algorithm checks for uniqueness in the code against any existing exams before finalizing it. The algorithm will retry generating codes with a new random seed if a collision exists in the existing room codes. The reduced ambiguous alphabet was decided after a pilot session when one of the students misread a capital letter O as 0 over the phone.

    \textbf{Deterministic Question Shuffle.} For each active student in the exam, questions are shuffled according to a pseudo random number sequence generated using the student ID and the exam ID. Thus, the shuffle order of questions remains the same every time a refresh is done by the same student but will be different for other active students taking the exam.

    \textbf{Mixed Scoring Pipeline}. MCQ items are auto-scored right away when a submission comes in. The auto-score is saved right away. Written and coding items' statuses are changed to "Evaluation Pending". After saving the manual evaluation by the teacher, the manual score is saved. After this, the total score is calculated as the sum of auto and manual scores. The submission's evaluation status is set to "Fully Evaluated" if there is no "Pending" item left.

    \textbf{State Polling}. The waiting room periodically polls the backend every three seconds to determine the time the exam starts. Teacher participant panel periodically polls every three seconds for participant statuses. The proctoring panel periodically polls every three seconds to get new snapshots and events. Student exam control periodically polls every two seconds so that freezing/force-submitting actions work instantly after pressing the corresponding button on the teacher panel's side. The Electron focus enforcer wakes up every five hundred milliseconds. The forbidden processes scaner scans the process list every one second.

    \textbf{Code Execution.} For each run request, the code is written into a temporary file under a unique per-request directory. In case of JavaScript files, code execution is done by the Node.js interpreter installed on the host machine. For Python files, execution is done by the python interpreter provided by the system. The c++ program is first compiled with g++ using a time-out limit of twelve seconds, after which it is executed within a limit of eight seconds. 

    \textbf{Event Logic for Proctoring.} The blaze face runs in the renderer at a low frame rate providing information about face detection count and pose information. The event class logic is generated based on the above-mentioned parameters; no face detected event, multiple faces detected event, looking away event, and looking down event. Events are rate-limited so that a particular event does not occur multiple times within a few seconds and every three seconds a snapshot is taken when the camera is active.

    \textbf{Desktop Integrity Control.} The Electron main process keeps the exam window in fullscreen and always on top during an active exam. It blocks shortcuts that would minimise or close the window, listens for fullscreen exit and focus loss events, scans the running process list against a forbidden list, and writes structured violation events to the backend with a severity tag.
  }
\end{justify}

%==============================================================
\subsection{Sequence Diagrams}

\begin{justify}
  {\fontsize{12}{18}\selectfont
    The following three sequence diagrams cover the main interaction flows for each user role. Figure~\ref{fig:3.8} shows the Teacher exam creation flow, Figure~\ref{fig:3.9} the Student exam attempt flow, and Figure~\ref{fig:3.10} the Admin user management flow.
  }
\end{justify}

%--- Figure 3.8: Sequence – Teacher Flow ---
\begin{figure}[t]
\centering
\begin{tikzpicture}[x=1cm,y=0.72cm,
  lifeline/.style={draw,dashed,gray!50,line width=0.8pt},
  msg/.code={\pgfkeysalso{->,>=Stealth,line width=0.9pt,draw=#1}},
  actor/.code={\pgfkeysalso{draw=#1,fill=#1,fill opacity=0.15,rounded corners=3pt,
                minimum width=2.2cm,minimum height=0.6cm,
                align=center,font=\scriptsize\bfseries}},
  note/.style={font=\tiny,anchor=west}
]
  % ---- Lifeline headers ----
  \node[actor=green!60!black] (tui)  at (0,0)  {Teacher\\UI};
  \node[actor=blue!60]        (api)  at (3.5,0) {Express\\API};
  \node[actor=orange!60]      (db)   at (7,0)   {PostgreSQL};

  % ---- Lifelines ----
  \draw[lifeline] (0,-0.4)  -- (0,-10.5);
  \draw[lifeline] (3.5,-0.4)-- (3.5,-10.5);
  \draw[lifeline] (7,-0.4)  -- (7,-10.5);

  % ---- Messages ----
  \foreach \y/\col/\fromx/\tox/\lbl in {
    -1.0/green!60!black/0/3.5/{POST /api/auth/login},
    -1.7/blue!60/3.5/0/{200 + JWT token},
    -2.6/green!60!black/0/3.5/{POST /api/exams (title, duration, type)},
    -3.3/blue!60/3.5/0/{201 exam\_id + room\_code},
    -4.2/green!60!black/0/3.5/{POST /api/exams/:id/questions (×N)},
    -4.9/blue!60/3.5/0/{201 question created},
    -5.8/green!60!black/0/3.5/{POST /api/exams/:id/start},
    -6.5/blue!60/3.5/0/{200 started\_at},
    -7.4/green!60!black/0/3.5/{GET /api/exams/:id/participants (poll)},
    -8.1/blue!60/3.5/0/{200 participant list + violations},
    -9.0/green!60!black/0/3.5/{PUT participants/:id/toggle-freeze},
    -9.7/blue!60/3.5/0/{200 frozen}
  }{
    \draw[msg=\col] (\fromx,\y) -- (\tox,\y)
        node[midway,above,font=\tiny]{\lbl};
  }
  % DB interactions
  \draw[msg=blue!60] (3.5,-1.0) -- (7,-1.0)
      node[midway,above,font=\tiny]{SELECT users};
  \draw[msg=blue!60] (7,-1.7) -- (3.5,-1.7)
      node[midway,above,font=\tiny]{user row};
  \draw[msg=blue!60] (3.5,-2.6) -- (7,-2.6)
      node[midway,above,font=\tiny]{INSERT exams};
  \draw[msg=blue!60] (7,-3.3) -- (3.5,-3.3)
      node[midway,above,font=\tiny]{exam record};

  % Activation boxes
  \fill[green!20,opacity=0.4] (-0.12,-0.8) rectangle (0.12,-10.2);
  \fill[blue!20, opacity=0.4] (3.38,-0.8) rectangle (3.62,-10.2);
  \fill[orange!20,opacity=0.4] (6.88,-0.8) rectangle (7.12,-4.5);
\end{tikzpicture}
\caption{Sequence diagram: Teacher exam creation and monitoring flow.}
\label{fig:3.8}
\end{figure}

%--- Figure 3.9: Sequence – Student Flow ---
\begin{figure}[t]
\centering
\begin{tikzpicture}[x=1cm,y=0.72cm,
  lifeline/.style={draw,dashed,gray!50,line width=0.8pt},
  msg/.code={\pgfkeysalso{->,>=Stealth,line width=0.9pt,draw=#1}},
  actor/.code={\pgfkeysalso{draw=#1,fill=#1,fill opacity=0.15,rounded corners=3pt,
                minimum width=2.0cm,minimum height=0.6cm,
                align=center,font=\scriptsize\bfseries}}
]
  % Lifeline headers
  \node[actor=blue!60]      (sui)  at (0,0)  {Student\\UI};
  \node[actor=gray!60]      (elec) at (3.2,0) {Electron\\Main};
  \node[actor=green!60]     (api)  at (6.4,0) {Express\\API};
  \node[actor=orange!60]    (db)   at (9.6,0) {PostgreSQL};

  \draw[dashed,gray!40] (0,-0.4)   -- (0,-12.0);
  \draw[dashed,gray!40] (3.2,-0.4) -- (3.2,-12.0);
  \draw[dashed,gray!40] (6.4,-0.4) -- (6.4,-12.0);
  \draw[dashed,gray!40] (9.6,-0.4) -- (9.6,-12.0);

  % Messages
  \draw[msg=blue!60]  (0,-1.0)  -- (6.4,-1.0) node[midway,above,font=\tiny]{POST /exams/join-by-room (roomCode, roll)};
  \draw[msg=orange!60](6.4,-1.6)-- (9.6,-1.6) node[midway,above,font=\tiny]{SELECT student + exam};
  \draw[msg=blue!60]  (6.4,-2.2)-- (0,-2.2)   node[midway,above,font=\tiny]{200 + JWT token + exam info};
  \draw[msg=blue!60]  (0,-3.2)  -- (6.4,-3.2) node[midway,above,font=\tiny]{GET /exams/:id/status (poll every 3 s)};
  \draw[msg=blue!60]  (6.4,-3.8)-- (0,-3.8)   node[midway,above,font=\tiny]{waiting};
  \draw[msg=blue!60]  (6.4,-4.8)-- (0,-4.8)   node[midway,above,font=\tiny]{in\_progress → exam starts};
  \draw[msg=blue!60]  (0,-5.5)  -- (3.2,-5.5) node[midway,above,font=\tiny]{ipcRenderer: start-exam};
  \draw[msg=gray!60]  (3.2,-6.1)-- (0,-6.1)   node[midway,above,font=\tiny]{fullscreen + lock shortcuts};
  \draw[msg=blue!60]  (0,-7.1)  -- (6.4,-7.1) node[midway,above,font=\tiny]{POST /exams/:id/run-code (code, lang)};
  \draw[msg=blue!60]  (6.4,-7.7)-- (0,-7.7)   node[midway,above,font=\tiny]{200 stdout / stderr};
  \draw[msg=blue!60]  (0,-8.7)  -- (6.4,-8.7) node[midway,above,font=\tiny]{POST /proctoring/:id/snapshot (base64)};
  \draw[msg=blue!60]  (6.4,-9.3)-- (0,-9.3)   node[midway,above,font=\tiny]{200 ok};
  \draw[msg=blue!60]  (0,-10.3) -- (6.4,-10.3)node[midway,above,font=\tiny]{POST /exams/:id/submit (answers)};
  \draw[msg=orange!60](6.4,-10.9)-- (9.6,-10.9)node[midway,above,font=\tiny]{INSERT submissions};
  \draw[msg=blue!60]  (6.4,-11.5)-- (0,-11.5) node[midway,above,font=\tiny]{201 submission\_id};

  \fill[blue!15, opacity=0.4] (-0.12,-0.8) rectangle (0.12,-11.8);
  \fill[gray!15, opacity=0.4] (3.08,-5.3)  rectangle (3.32,-6.5);
\end{tikzpicture}
\caption{Sequence diagram: Student exam attempt flow from room entry to submission.}
\label{fig:3.9}
\end{figure}

%--- Figure 3.10: Sequence – Admin Flow ---
\begin{figure}[t]
\centering
\begin{tikzpicture}[x=1cm,y=0.72cm,
  lifeline/.style={draw,dashed,gray!50,line width=0.8pt},
  msg/.code={\pgfkeysalso{->,>=Stealth,line width=0.9pt,draw=#1}},
  actor/.code={\pgfkeysalso{draw=#1,fill=#1,fill opacity=0.15,rounded corners=3pt,
                minimum width=2.2cm,minimum height=0.6cm,
                align=center,font=\scriptsize\bfseries}}
]
  \node[actor=violet!60]  (aui) at (0,0)   {Admin\\UI};
  \node[actor=green!60]   (api) at (3.8,0) {Express\\API};
  \node[actor=orange!60]  (db)  at (7.6,0) {PostgreSQL};

  \draw[dashed,gray!40] (0,-0.4)   -- (0,-10.0);
  \draw[dashed,gray!40] (3.8,-0.4) -- (3.8,-10.0);
  \draw[dashed,gray!40] (7.6,-0.4) -- (7.6,-10.0);

  \draw[msg=violet!60] (0,-1.0)  -- (3.8,-1.0) node[midway,above,font=\tiny]{POST /auth/login (admin credentials)};
  \draw[msg=green!60]  (3.8,-1.6)-- (7.6,-1.6) node[midway,above,font=\tiny]{SELECT user WHERE role=admin};
  \draw[msg=violet!60] (3.8,-2.2)-- (0,-2.2)   node[midway,above,font=\tiny]{200 + JWT token};

  \draw[msg=violet!60] (0,-3.2)  -- (3.8,-3.2) node[midway,above,font=\tiny]{GET /admin/stats};
  \draw[msg=green!60]  (3.8,-3.8)-- (7.6,-3.8) node[midway,above,font=\tiny]{COUNT users, exams};
  \draw[msg=violet!60] (3.8,-4.4)-- (0,-4.4)   node[midway,above,font=\tiny]{200 \{teachers, students, exams\}};

  \draw[msg=violet!60] (0,-5.4)  -- (3.8,-5.4) node[midway,above,font=\tiny]{POST /admin/create-teacher \{email, password\}};
  \draw[msg=green!60]  (3.8,-6.0)-- (7.6,-6.0) node[midway,above,font=\tiny]{INSERT users (role=teacher)};
  \draw[msg=violet!60] (3.8,-6.6)-- (0,-6.6)   node[midway,above,font=\tiny]{201 user record};

  \draw[msg=violet!60] (0,-7.6)  -- (3.8,-7.6) node[midway,above,font=\tiny]{POST /admin/upload-students \{csvContent\}};
  \draw[msg=green!60]  (3.8,-8.2)-- (7.6,-8.2) node[midway,above,font=\tiny]{bulk INSERT students};
  \draw[msg=violet!60] (3.8,-8.8)-- (0,-8.8)   node[midway,above,font=\tiny]{200 \{created: N, skipped: M\}};

  \draw[msg=violet!60] (0,-9.6)  -- (3.8,-9.6) node[midway,above,font=\tiny]{PUT /admin/users/:id/status \{status:inactive\}};
  \draw[msg=violet!60] (3.8,-9.6)-- (0,-9.6)   {};

  \fill[violet!15,opacity=0.4] (-0.12,-0.8) rectangle (0.12,-9.9);
\end{tikzpicture}
\caption{Sequence diagram: Admin user management flow including bulk student upload.}
\label{fig:3.10}
\end{figure}

%==============================================================
\subsection{Data Flow Diagrams}

\begin{justify}
  {\fontsize{12}{18}\selectfont
    Figure~\ref{fig:3.3} shows the Level 0 DFD, treating the entire system as a single process. Figure~\ref{fig:3.4} expands this into the Level 1 DFD with six internal processes and five data stores.
  }
\end{justify}

%--- Figure 3.3: DFD Level 0 ---
\begin{figure}[t]
\centering
\begin{tikzpicture}[node distance=0pt]
  % Central process
  \node[draw=green!55!black,fill=green!8,circle,
        minimum size=3.4cm,align=center,font=\small\bfseries,
        text=green!40!black] (p0) at (0,0)
        {P0\\Invigilo\\Exam System};

  % External entities
  \node[entity] (student)  at (-5.5, 0)   {Student};
  \node[entity] (teacher)  at ( 5.5, 0)   {Teacher};
  \node[entity] (admin)    at ( 0,   3.6)  {Admin};

  % Student flows
  \draw[arr] (student.east) -- ([xshift=-1.7cm]p0.west)
      node[midway,above,font=\scriptsize]{Credentials, Answers, Proct.\ Data};
  \draw[arr] ([xshift=-1.7cm]p0.west) -- (student.east |- 0,-0.45)
      node[midway,below,font=\scriptsize]{Auth Status, Exam State, Results};

  % Teacher flows
  \draw[arr] (teacher.west) -- ([xshift=1.7cm]p0.east)
      node[midway,above,font=\scriptsize]{Exam Setup, Evaluation};
  \draw[arr] ([xshift=1.7cm]p0.east) -- (teacher.west |- 0,-0.45)
      node[midway,below,font=\scriptsize]{Participant Status, Proct.\ View};

  % Admin flows
  \draw[arr] (admin.south) -- (p0.north)
      node[midway,right,font=\scriptsize]{User Management Commands};
  \draw[arr] ([xshift=4pt]p0.north) -- ([xshift=4pt]admin.south)
      node[midway,left,font=\scriptsize]{User Lists, Stats};
\end{tikzpicture}
\caption{DFD Level 0: Invigilo system as a single process with external entities.}
\label{fig:3.3}
\end{figure}

%--- Figure 3.4: DFD Level 1 ---
\begin{figure}[t]
\centering
\begin{tikzpicture}[x=1cm,y=1cm]
  % ---- External entities (left column, x=0) ----
  \node[entity,minimum width=1.6cm] (adm)  at (0, 5.5) {Admin};
  \node[entity,minimum width=1.6cm] (tch1) at (0, 3.5) {Teacher};
  \node[entity,minimum width=1.6cm] (stu1) at (0, 1.5) {Student};
  \node[entity,minimum width=1.6cm] (tch2) at (0,-0.5) {Teacher};
  \node[entity,minimum width=1.6cm] (stu2) at (0,-2.5) {Student};

  % ---- Processes (centre column, x=5) ----
  \node[process] (p6) at (5, 5.5) {\small P6\\User\\Admin};
  \node[process] (p2) at (5, 3.5) {\small P2\\Exam\\Setup};
  \node[process] (p1) at (5, 1.5) {\small P1\\Auth};
  \node[process] (p3) at (5,-0.5) {\small P3\\Exam\\Attempt};
  \node[process] (p5) at (5,-2.5) {\small P5\\Evaluation\\+ Proctor};

  % ---- Data stores (right column, x=9.5) ----
  \node[dstore,minimum width=2.6cm] (d1) at (9.5, 5.5) {D1: Users};
  \node[dstore,minimum width=2.6cm] (d2) at (9.5, 3.5) {D2: Exams \&\\Questions};
  \node[dstore,minimum width=2.6cm] (d3) at (9.5, 1.5) {D3: Participants};
  \node[dstore,minimum width=2.6cm] (d4) at (9.5,-0.5) {D4: Submissions};
  \node[dstore,minimum width=2.6cm] (d5) at (9.5,-2.5) {D5: Proct.\ Data};

  % ---- Entity → Process arrows ----
  \draw[arr] (adm)  -- (p6);
  \draw[arr] (tch1) -- (p2);
  \draw[arr] (stu1) -- (p1);
  \draw[arr] (tch2) -- (p3);
  \draw[arr] (stu2) -- (p5);

  % ---- Process → Data store arrows ----
  \draw[arr] (p6) -- (d1);
  \draw[arr] (p2) -- (d2);
  \draw[darr] (p1) -- (d3) node[midway,above,font=\tiny]{verify/join};
  \draw[arr] (p3) -- (d4);
  \draw[darr] (p5) -- (d5);

  % ---- Cross-process data reads (avoid crossing: use bent paths) ----
  % P1 reads D1 (auth lookup) — bend above
  \draw[arr,dashed,gray!60]
      (d1.south) to[out=270,in=0] (p1.east)
      node[near start,right,font=\tiny]{lookup};
  % P3 reads D2 (questions) — straight right then down
  \draw[arr,dashed,gray!60]
      (d2.south) to[out=270,in=0] (p3.east)
      node[near start,right,font=\tiny]{questions};
  % P5 reads D4 (submissions for eval)
  \draw[darr,dashed,gray!60]
      (d4.south) to[out=270,in=0] (p5.east)
      node[near start,right,font=\tiny]{eval};

  % ---- Process → Entity output labels ----
  \draw[arr] (p1) to[out=180,in=0] ([yshift=6pt]stu1.east)
      node[midway,above,font=\tiny]{token};
  \draw[arr] (p5) to[out=180,in=0] ([yshift=6pt]stu2.east)
      node[midway,above,font=\tiny]{results};
\end{tikzpicture}
\caption{DFD Level 1: Six internal processes and five data stores of Invigilo.}
\label{fig:3.4}
\end{figure}

%==============================================================
\subsection{Entity Relationship Diagram}

\begin{justify}
  {\fontsize{12}{18}\selectfont
    Figure~\ref{fig:3.5} presents the ER diagram of the seven core tables in the Invigilo PostgreSQL database.
  }
\end{justify}

%--- Figure 3.5: ER Diagram ---
\begin{figure}[t]
\centering
\begin{tikzpicture}[
  ent/.style={draw=blue!60,fill=blue!8,rounded corners=5pt,
              align=left,font=\scriptsize,inner sep=5pt,
              minimum width=3.0cm},
  rel/.style={font=\tiny,draw=gray!50,fill=gray!10,
              rounded corners=3pt,inner sep=2pt},
  earr/.style={arr,gray!55,line width=1pt}
]
  % ---- Entities ----
  \node[ent] (users) at (0,0) {
    \textbf{USERS}\\id (PK)\\name, email (UK)\\password\_hash\\roll\_number, role\\status, created\_at
  };

  \node[ent] (exams) at (5.5,2.5) {
    \textbf{EXAMS}\\id (PK)\\title, exam\_type\\duration, status\\room\_code (UK)\\created\_by (FK)\\webcam\_required
  };

  \node[ent] (questions) at (11,2.5) {
    \textbf{QUESTIONS}\\id (PK)\\exam\_id (FK)\\question\_text\\question\_type\\options (JSONB)\\correct\_answer, marks
  };

  \node[ent] (submissions) at (11,-2.0) {
    \textbf{SUBMISSIONS}\\id (PK)\\exam\_id (FK)\\student\_id (FK)\\answers (JSONB)\\auto\_score, manual\_score\\evaluation\_status
  };

  \node[ent] (participants) at (5.5,-2.0) {
    \textbf{EXAM\_PARTICIPANTS}\\id (PK)\\exam\_id (FK)\\student\_id (FK)\\status, violation\_count\\is\_frozen
  };

  \node[ent] (pevents) at (0,4.2) {
    \textbf{PROCT.\_EVENTS}\\id (PK)\\exam\_id (FK)\\student\_id (FK)\\event\_type, details
  };

  \node[ent] (psnaps) at (0,-3.5) {
    \textbf{PROCT.\_SNAPSHOTS}\\exam\_id (PK,FK)\\student\_id (PK,FK)\\snapshot\_base64\\face\_count, status
  };

  % ---- Relationships ----
  % USERS creates EXAMS
  \draw[earr] (users.east |- exams.west) -- (exams.west)
      node[midway,above,rel]{creates};
  % EXAMS contains QUESTIONS
  \draw[earr] (exams.east) -- (questions.west)
      node[midway,above,rel]{contains};
  % USERS joins EXAM\_PARTICIPANTS
  \draw[earr] (users.south) to[out=270,in=180] (participants.west)
      node[near end,above,rel]{joins};
  % EXAMS has EXAM\_PARTICIPANTS
  \draw[earr] (exams.south) -- (participants.north)
      node[midway,right,rel]{has};
  % USERS submits SUBMISSIONS
  \draw[earr] (users.east |- submissions.west)
      to[out=0,in=180] (submissions.west)
      node[near end,above,rel]{submits};
  % EXAMS receives SUBMISSIONS
  \draw[earr] (exams.east |- submissions.north)
      to[out=0,in=90] (submissions.north)
      node[near start,right,rel]{receives};
  % USERS triggers PROCT\_EVENTS
  \draw[earr] (users.north) -- (pevents.south)
      node[midway,right,rel]{triggers};
  % USERS uploads PROCT\_SNAPSHOTS
  \draw[earr] (users.south) to[out=260,in=90] (psnaps.north)
      node[near end,left,rel]{uploads};
\end{tikzpicture}
\caption{Entity relationship diagram of the Invigilo database schema.}
\label{fig:3.5}
\end{figure}

%==============================================================
\subsection{Activity Diagram}

\begin{justify}
  {\fontsize{12}{18}\selectfont
    Figure~\ref{fig:3.6} shows the activity diagram for the full exam lifecycle, capturing the parallel activities of the teacher and the student from exam creation through evaluation.
  }
\end{justify}

%--- Figure 3.6: Activity Diagram ---
\begin{figure}[t]
\centering
\begin{tikzpicture}[
  act/.code={\pgfkeysalso{draw=#1,fill=#1,fill opacity=0.12,rounded corners=6pt,
              minimum width=3.4cm,minimum height=0.75cm,
              align=center,font=\scriptsize}},
  decd/.style={draw=orange!70,fill=orange!10,diamond,
               minimum width=2.6cm,minimum height=0.75cm,
               align=center,font=\scriptsize,aspect=3.0},
  fa/.style={arr,gray!55,line width=1pt},
  sync/.style={draw,fill=black,minimum width=6.0cm,minimum height=0.12cm}
]
  % ---- Swimlane headers ----
  \draw[gray!30,fill=green!5]  (-0.3,1.0) rectangle (3.8,-14.2);
  \draw[gray!30,fill=blue!5]   (4.2,1.0)  rectangle (8.3,-14.2);
  \node[font=\small\bfseries,text=green!55!black] at (1.75,0.65)  {Teacher};
  \node[font=\small\bfseries,text=blue!65!black]  at (6.25,0.65)  {Student};
  \draw[gray!40] (4.0,1.0) -- (4.0,-14.2);

  % ---- START ----
  \node[draw,fill=black,circle,minimum size=0.4cm] (start) at (4.0,0.2) {};

  % ---- Teacher side ----
  \node[act=green!55!black] (tcreate) at (1.75,-0.9)  {Create Exam \& Add Questions};
  \node[act=green!55!black] (tshare)  at (1.75,-2.0)  {Share Room Code};
  \node[act=green!55!black] (tmonit)  at (1.75,-4.5)  {Monitor Participants};
  \node[decd]               (tact)    at (1.75,-5.7)  {Action?};
  \node[act=red!60]         (tfreeze) at (1.75,-7.0)  {Freeze / Force Submit};
  \node[act=green!55!black] (teval)   at (1.75,-9.5)  {Evaluate Written \& Coding};
  \node[act=green!55!black] (tfinish) at (1.75,-10.7) {Publish Scores};

  % ---- Student side ----
  \node[act=blue!65!black]  (sjoin)   at (6.25,-2.0)  {Join Waiting Room};
  \node[act=blue!65!black]  (swait)   at (6.25,-3.2)  {Wait for Exam Start};
  \node[act=blue!65!black]  (sexam)   at (6.25,-4.5)  {Attempt Questions};
  \node[act=blue!65!black]  (srun)    at (6.25,-5.7)  {Run Code (if coding Q)};
  \node[decd]               (ssub)    at (6.25,-7.0)  {Submit?};
  \node[act=blue!65!black]  (sdone)   at (6.25,-8.2)  {Submission Recorded};
  \node[act=blue!65!black]  (sresult) at (6.25,-10.7) {View Detailed Results};

  % ---- Sync bars ----
  \node[sync] (sync1) at (4.0,-1.45) {};   % After create, before join
  \node[sync] (sync2) at (4.0,-3.75) {};   % Exam starts
  \node[sync] (sync3) at (4.0,-8.85) {};   % Submission done, eval begins

  % ---- END ----
  \node[draw,fill=black,circle,minimum size=0.4cm,
        path picture={\draw[white,line width=1.5pt]
        (path picture bounding box.center) circle (0.12cm);}]
       (endn) at (4.0,-11.8) {};

  % ---- Flow lines ----
  % From start
  \draw[fa] (start) -- (sync1);
  % Teacher side
  \draw[fa] (sync1.west|-tcreate) -- (tcreate);
  \draw[fa] (tcreate) -- (tshare);
  \draw[fa] (tshare.south|-sync1.south) -- (sync1.west|-tshare) {};
  % Student side
  \draw[fa] (sync1.east|-sjoin) -- (sjoin);
  % Exam start sync
  \draw[fa] (tshare) -- (tshare |- sync2) -- (sync2);
  \draw[fa] (sjoin) -- (sjoin |- sync2) -- (sync2);
  \draw[fa] (swait.north) -- (sync2.east|-swait);

  \draw[fa] (sync2) |- (tmonit);
  \draw[fa] (sync2) |- (swait);
  \draw[fa] (swait) -- (sexam);

  \draw[fa] (tmonit) -- (tact);
  \draw[fa] (tact.south) -- (tfreeze) node[midway,left,font=\tiny]{violation};
  \draw[fa] (tact.east) to[out=0,in=0] (tmonit.east)
      node[near start,right,font=\tiny]{observe};

  \draw[fa] (sexam) -- (srun);
  \draw[fa] (srun) -- (ssub);
  \draw[fa] (ssub.south) -- (sdone) node[midway,left,font=\tiny]{yes};
  \draw[fa] (ssub.east) to[out=0,in=0] (sexam.east)
      node[near start,right,font=\tiny]{no};

  % Submission sync
  \draw[fa] (sdone) -- (sdone |- sync3) -- (sync3);
  \draw[fa] (tfreeze |- sync3) -- (sync3);

  \draw[fa] (sync3) |- (teval);
  \draw[fa] (teval) -- (tfinish);
  \draw[fa] (sync3) |- (sresult);

  % End
  \draw[fa] (tfinish.south|-endn) -- (endn);
  \draw[fa] (sresult.south|-endn) -- (endn);
\end{tikzpicture}
\caption{Activity diagram showing parallel Teacher and Student activities across the exam lifecycle.}
\label{fig:3.6}
\end{figure}

%==============================================================
\subsection{Class Diagram}

\begin{justify}
  {\fontsize{12}{18}\selectfont
    Figure~\ref{fig:3.7} presents the class diagram for the core backend controllers and their relationships to the domain entities.
  }
\end{justify}

%--- Figure 3.7: Class Diagram ---
\begin{figure}[t]
\centering
\begin{tikzpicture}[
  cls/.style={draw=blue!50,fill=blue!6,rounded corners=4pt,
              font=\scriptsize,inner sep=0pt,minimum width=3.0cm},
  clsh/.style={draw=blue!50,fill=blue!20,rounded corners=4pt,
               font=\small\bfseries,inner sep=4pt,
               minimum width=3.0cm,align=center},
  assoc/.style={->,>=Stealth,gray!60,line width=0.9pt},
  dep/.style={->,>=Stealth,gray!50,dashed,line width=0.8pt}
]
  % Helper macro: draw a UML class box
  % \umlclass{name}{x}{y}{title}{attributes}{methods}
  % We draw as nested nodes

  % ---- Domain entity classes (top row) ----
  \foreach \nm/\cx/\ttl/\attrs in {
    User/0/User/{+id: int\\+name, email: str\\+role, status: str},
    Exam/4.5/Exam/{+id: int\\+title, exam\_type: str\\+duration: int\\+status, room\_code: str},
    Question/9/Question/{+id: int\\+exam\_id: int\\+question\_type: str\\+marks: int},
    Submission/13.5/Submission/{+id: int\\+answers: JSON\\+auto\_score: int\\+evaluation\_status: str}
  }{
    \node[clsh] (\nm H) at (\cx, 0.4)  {\ttl};
    \node[cls,anchor=north,text width=2.8cm,align=left,
          minimum height=1.5cm] (\nm) at (\cx, 0.0) {\attrs};
  }

  % ---- Controller classes (bottom row) ----
  \foreach \nm/\cx/\ttl/\meths in {
    AuthCtrl/0/AuthController/{+login()\\+getCurrentUser()},
    ExamCtrl/4.5/ExamController/{+createExam()\\+joinExam()\\+startExam()\\+submitExam()\\+forceSubmit()},
    QCtrl/9/QuestionController/{+addQuestion()\\+evaluateAnswers()},
    AdminCtrl/13.5/AdminController/{+createTeacher()\\+createStudent()\\+uploadStudents()\\+toggleStatus()}
  }{
    \node[clsh] (\nm H) at (\cx,-4.8) {\ttl};
    \node[cls,anchor=north,text width=2.8cm,align=left,
          minimum height=1.6cm] (\nm) at (\cx,-5.2) {\meths};
  }

  % Additional controllers
  \node[clsh] (ProcH)  at (4.5,-8.6) {ProctoringController};
  \node[cls,anchor=north,text width=2.8cm,align=left,
        minimum height=1.4cm] (Proc) at (4.5,-9.0)
       {+reportEvent()\\+uploadSnapshot()\\+getStudentsStatus()};

  \node[clsh] (AIH)    at (9,-8.6)   {AIController};
  \node[cls,anchor=north,text width=2.8cm,align=left,
        minimum height=1.1cm] (AI) at (9,-9.0)
       {+chat()\\+generateQuestions()};

  % ---- Associations (entity relationships) ----
  \draw[assoc] (UserH.east) -- (ExamH.west)  node[midway,above,font=\tiny]{creates};
  \draw[assoc] (ExamH.east) -- (QuestionH.west) node[midway,above,font=\tiny]{contains};
  \draw[assoc] (UserH.south) to[out=270,in=270,looseness=0.4]
               (SubmissionH.south) node[near end,below,font=\tiny]{submits};

  % ---- Dependencies (controllers use entities) ----
  \draw[dep] (AuthCtrl) -- (User);
  \draw[dep] (ExamCtrl) -- (Exam);
  \draw[dep] (QCtrl) -- (Question);
  \draw[dep] (AdminCtrl) -- (User);
  \draw[dep] (Proc) to[out=90,in=270] (ExamH.south);
  \draw[dep] (AI) to[out=90,in=270] (ExamH.south);
\end{tikzpicture}
\caption{Class diagram of the Invigilo backend controllers and domain entities.}
\label{fig:3.7}
\end{figure}

\section{Conclusion}

\begin{justify}
  {\fontsize{12}{18}\selectfont
    This methodology is intentionally practical. Everything discussed in the algorithms and workflows sections was derived from what had to be built, tested, and demonstrated. The distinct boundaries between the Electron main process, the React rendering process, the Express server, and the PostgreSQL database enabled the development team to change each layer independently without affecting other layers. The invariants at the workflow level were selected for their simplicity to understand and validate manually in the live demo context. Chapter 4 illustrates how this methodology was implemented and provides results obtained from the implementation.
  }
\end{justify}

%==============================================================
% CHAPTER IV: Implementation, Results and Discussions
%==============================================================
\chapter{Implementation, Results and Discussions}

\section{Introduction}

\begin{justify}
  {\fontsize{12}{18}\selectfont
    This chapter reports on the implementation of Invigilo as it stands in the develop branch of the repository as of April 2026, the experimental setup used for verification, the metrics collected from the running system, and a manual integration test suite that exercises the full exam lifecycle. The chapter ends with a mapping from the original objectives stated in Chapter 1 to the evidence available in the current build.
  }
\end{justify}

\section{Experimental Setup}

\begin{justify}
  {\fontsize{12}{18}\selectfont
    Verification was done on Windows 11 host computers using Node.js version 20, PostgreSQL version 16, and Electron 39.2.7. The frontend was developed using Vite version 7.1.7. The React 19.2.0 framework and Tailwind CSS were used to build the renderer. Express 4.18.2 was used in the backend, jsonwebtoken was used for authentication tokens generation, bcryptjs for hashing passwords, and the pg driver was used for access to PostgreSQL database. The BlazeFace was run using TensorFlow.js in the renderer while model assets were served from the static folder of the renderer.

    End-to-end verification was conducted using two laptops. The first one was used to host the backend server, the database, and the dashboard for the teacher. The other one was utilized for running the desktop client and taking the test itself. Wireless local area network interconnected both machines. Such approach realistically simulates the laboratory quiz environment which will help verify the functionality of polling, freezing, and forced submission as well as snapshotting pipelines by the proctor.

   The build verification step ran the build of the renderer using \texttt{npm run build:renderer} command. During several attempts the build passed successfully with exit code zero and execution times from eleven to thirteen seconds; specifically, Vite build engine reported ten-and-a-half-second-long internal build process. No build errors or warnings with a greater severity than info were encountered during this verification step.

   While performance was not in the spotlight during integration testing, we made sure that the following features function correctly: roles change workflow, API contract consistency between renderer and backend, submission and evaluation pipeline, desktop lock mechanism, and erroneous credential handling.
  }
\end{justify}

\section{Evaluation Metrics}

\begin{justify}
  {\fontsize{12}{18}\selectfont
    Because Invigilo is an applied system rather than a research artifact, the evaluation metrics used are mostly system level and behavioural, not statistical. The metrics are listed below.

    \begin{itemize}
        \item \textbf{Functional Coverage:} number of role specific capabilities implemented and verified end to end.
        \item \textbf{API Surface:} number of REST endpoints exposed by the backend, broken down by module.
        \item \textbf{Schema Surface:} number of relational tables used to persist the system state.
        \item \textbf{Build Health:} success of the renderer production build and the wall clock build time.
        \item \textbf{Live Control Latency:} the maximum time between a teacher action (freeze, force submit) and the student observing the effect, derived from the configured polling intervals.
        \item \textbf{Test Pass Rate:} the fraction of manual integration test cases passing in the most recent verification run.
    \end{itemize}

    These metrics were chosen because they directly reflect whether the system fulfils its design objectives. A web based exam platform that builds successfully, exposes a clean API, persists structured data and reacts to teacher commands within a few seconds is the practical definition of a usable lab quiz tool.
  }
\end{justify}

\section{Dataset / System Parameters}

\begin{justify}
  {\fontsize{12}{18}\selectfont
    Invigilo is not trained on a dataset in the machine learning sense. The face detection model used internally by BlazeFace was pretrained by its original authors and shipped through TensorFlow.js. For the purpose of this report the relevant equivalent of a dataset is the set of internal system parameters that govern timing, control and resource limits. Table 4.1 lists these parameters and the values used in the verified build.
  }
\end{justify}

\begin{table}[H]
  \caption{Runtime Control and Monitoring Parameters}
  \label{tab:4.1}
  \centering
  {\fontsize{11}{14}\selectfont
    \begin{tabular}{|p{2.6in}|c|}
      \hline
      \textbf{Parameter} & \textbf{Value} \\
      \hline
      Waiting room status poll interval & 3 seconds \\
      \hline
      Teacher participant poll interval & 3 seconds \\
      \hline
      Teacher proctoring poll interval & 3 seconds \\
      \hline
      Student exam control poll interval & 2 seconds \\
      \hline
      Webcam snapshot upload interval & 3 seconds \\
      \hline
      Electron focus enforcer interval & 0.5 seconds \\
      \hline
      Forbidden process scan interval & 1 second \\
      \hline
      Warning display timeout & 4 seconds \\
      \hline
      Code execution timeout & 8 seconds \\
      \hline
      C++ compilation timeout & 12 seconds \\
      \hline
      Maximum allowed violations before auto submit & 3 \\
      \hline
      Room code length & 6 characters \\
      \hline
    \end{tabular}
  }
\end{table}

\section{Implementation and Results}

\subsection{Qualitative Results}

\begin{justify}
  {\fontsize{12}{18}\selectfont
    There are several qualitative observations that can be made about the process of conducting exams manually through integration:

    First and foremost, the teacher's dashboard made sure that everything concerning the exam was done within one window. As soon as the user logged in, the required options to create an exam, ask questions, share the room code, monitor participants' activities, freeze the session, make the participant complete the test, and work at the assessment table became available from the side panel of navigation. Nothing else but the same tool was used all the time during the session, and the teacher did not need more than one link address (the one for logging in).

    The participant panel got refreshed each time when the state changed for no more than three seconds. If the participant managed to log in, the student's counter would go up by one. If there was cheating, for instance, leaving the exam window or committing some other actions, the violation counter and the current type of violation would appear next to the participant's name. Freezing made an overlay block the participant's screen and interaction with the exam in two seconds.

    The experience of the student was limited in a sense that there was no idle time required on the part of the student because of how the examination procedure was carried out. There was no need for the student to click again and again the "Check Now" button because the system would automatically check for the availability of the server. Layout-wise, the examination page had the timer located at the top part of the page, a sidebar for the navigation of the question in all-at-once mode, and just a single question being presented in one-by-one mode. In case there were coding questions, a Monaco-based code editor will load with a code snippet loaded if there was any, the default choice of the programming language will be what the teacher set up, and there was standard output and standard error below the code editor after running the program.

    Low-ceremony in the administrative governance procedure means that the administrator will have to simply create a new account by filling out a simple form or even bulk-create 20 accounts by pasting a CSV fragment in the text area. Activation and deactivation of the accounts is done with just a single click. Dashboard metrics are updated immediately upon refreshing the dashboard.

    Finally, the security posture observed during the pilot was visibly stronger than what a normal browser based quiz would offer. Pressing Alt+F4 did not close the window. Pressing F11 did not toggle fullscreen. Trying to open Task Manager triggered a forbidden process violation that appeared in the teacher's dashboard within seconds. Attempting to start a screen sharing application also triggered a violation. None of this stops a determined attacker with administrative privileges, but it raises the cost of casual cheating well above the cost of just reading the question and answering it honestly.
  }
\end{justify}

%--- Figures 4.1–4.4: Application Screenshots ---
\begin{figure}[t]
\centering
\begin{tikzpicture}
  \draw[fill=gray!10,draw=gray!50,rounded corners=6pt]
      (0,0) rectangle (12.5,6.0);
  \node[font=\small\bfseries,text=gray!60] at (6.25,3.2) {Screenshot Placeholder};
  \node[font=\scriptsize,text=gray!50,align=center] at (6.25,2.5)
      {Teacher Dashboard --- Exam Management View\\
       (Replace with actual application screenshot)};
  \draw[gray!30,dashed] (0.4,0.4) rectangle (12.1,5.6);
\end{tikzpicture}
\caption{Teacher dashboard showing the exam management and question authoring interface.}
\label{fig:4.1}
\end{figure}

\begin{figure}[t]
\centering
\begin{tikzpicture}
  \draw[fill=gray!10,draw=gray!50,rounded corners=6pt]
      (0,0) rectangle (12.5,6.0);
  \node[font=\small\bfseries,text=gray!60] at (6.25,3.2) {Screenshot Placeholder};
  \node[font=\scriptsize,text=gray!50,align=center] at (6.25,2.5)
      {Student Exam Interface --- Question Attempt View\\
       (Replace with actual application screenshot)};
  \draw[gray!30,dashed] (0.4,0.4) rectangle (12.1,5.6);
\end{tikzpicture}
\caption{Student exam interface showing the question attempt screen with timer and question navigation.}
\label{fig:4.2}
\end{figure}

\begin{figure}[t]
\centering
\begin{tikzpicture}
  \draw[fill=gray!10,draw=gray!50,rounded corners=6pt]
      (0,0) rectangle (12.5,6.0);
  \node[font=\small\bfseries,text=gray!60] at (6.25,3.2) {Screenshot Placeholder};
  \node[font=\scriptsize,text=gray!50,align=center] at (6.25,2.5)
      {Teacher Live Proctoring View --- Participant Monitoring Grid\\
       (Replace with actual application screenshot)};
  \draw[gray!30,dashed] (0.4,0.4) rectangle (12.1,5.6);
\end{tikzpicture}
\caption{Teacher live proctoring panel showing webcam snapshots and violation counts for each participant.}
\label{fig:4.3}
\end{figure}

\begin{figure}[t]
\centering
\begin{tikzpicture}
  \draw[fill=gray!10,draw=gray!50,rounded corners=6pt]
      (0,0) rectangle (12.5,6.0);
  \node[font=\small\bfseries,text=gray!60] at (6.25,3.2) {Screenshot Placeholder};
  \node[font=\scriptsize,text=gray!50,align=center] at (6.25,2.5)
      {Admin Dashboard --- User Management and Statistics View\\
       (Replace with actual application screenshot)};
  \draw[gray!30,dashed] (0.4,0.4) rectangle (12.1,5.6);
\end{tikzpicture}
\caption{Admin dashboard showing platform statistics and user management interface.}
\label{fig:4.4}
\end{figure}

\subsection{Quantitative Results}

\begin{justify}
  {\fontsize{12}{18}\selectfont
    The quantitative numbers collected from the implementation are summarised in the tables below. Table 4.2 captures the scale of the system. Table 4.3 captures the result of the renderer production build. Table 4.4 captures the manual integration test pass rate.
  }
\end{justify}

\begin{table}[H]
  \caption{Project Scale Metrics}
  \label{tab:4.2}
  \centering
  {\fontsize{11}{14}\selectfont
    \begin{tabular}{|p{3in}|r|}
      \hline
      \textbf{Metric} & \textbf{Value} \\
      \hline
      Total backend API endpoints & 39 \\
      \hline
      Auth endpoints & 2 \\
      \hline
      Exam endpoints & 25 \\
      \hline
      Proctoring endpoints & 4 \\
      \hline
      AI assistant endpoints & 2 \\
      \hline
      Admin endpoints & 6 \\
      \hline
      Core relational tables & 7 \\
      \hline
      Core frontend lines of code (six major files) & 4991 \\
      \hline
      Core backend controller lines of code (seven files) & 3139 \\
      \hline
      Total Git commits across the project & 60 \\
      \hline
    \end{tabular}
  }
\end{table}

\begin{table}[H]
  \caption{Build Verification Snapshot}
  \label{tab:4.3}
  \centering
  {\fontsize{11}{14}\selectfont
    \begin{tabular}{|p{3in}|p{2in}|}
      \hline
      \textbf{Verification Item} & \textbf{Outcome} \\
      \hline
      Renderer production build & Passed \\
      \hline
      Build exit code & 0 \\
      \hline
      Vite reported build time & approximately 11.45 s \\
      \hline
      Measured wall clock build time & approximately 12.75 s \\
      \hline
      Build warnings of severity higher than info & none observed \\
      \hline
    \end{tabular}
  }
\end{table}

\begin{table}[H]
  \caption{Manual Integration Test Cases and Outcomes}
  \label{tab:4.4}
  \centering
  {\fontsize{10}{13}\selectfont
    \begin{tabular}{|p{0.6in}|p{1.6in}|p{2in}|p{0.7in}|}
      \hline
      \textbf{TC ID} & \textbf{Test Case} & \textbf{Expected Result} & \textbf{Status} \\
      \hline
      TC-01 & Exam creation with one by one mode and randomised order & Settings persist in teacher manage view and in student view & Pass \\
      \hline
      TC-02 & Waiting room auto check & Status auto refresh, no manual button needed & Pass \\
      \hline
      TC-03 & Teacher start exam transition & Exam moves to in progress, students enter exam screen & Pass \\
      \hline
      TC-04 & Live violation ingest & Teacher panel shows violation count and latest type & Pass \\
      \hline
      TC-05 & Freeze and unfreeze participant & Student input locks and unlocks within two seconds & Pass \\
      \hline
      TC-06 & Force submit participant & Student sees overlay, submission recorded & Pass \\
      \hline
      TC-07 & Coding answer rendering & Multiline code preserved on evaluation desk & Pass \\
      \hline
      TC-08 & Manual evaluation save & Manual score updates, total recalculates, status updated & Pass \\
      \hline
      TC-09 & Student result detail view & Per question marks and comments visible & Pass \\
      \hline
      TC-10 & Teacher delete permission guard & Non owner request rejected with permission error & Pass \\
      \hline
      TC-11 & Desktop lockdown and refocus & Fullscreen restored, shortcuts blocked, violations logged & Pass \\
      \hline
      TC-12 & Renderer production build & Build completes with exit code zero & Pass \\
      \hline
    \end{tabular}
  }
\end{table}

\subsection{Analysis of the Results}

\begin{justify}
  {\fontsize{12}{18}\selectfont
    The numbers above reflect that there is sufficient scalability in the implemented system in terms of being used at a department-level system. Thirty nine data points exist within seven modules to cover all processes of creating an account through viewing results. Seven relational database tables are used to normalise users, exams, questions, participants, submissions, answers and violations. The system renders in less than eleven seconds, which is efficient enough in terms of the developer's ergonomic considerations.

    What should be mentioned next is the polling frequencies chosen for the live system. They guarantee that almost each change of state will be perceived with less than three seconds delay, while freeze and force-submit operations have a maximum delay of less than two seconds. In such a way, it is possible to respond in time to a student who misbehaves in class, while not flooding the system with too many requests.

    Twelve out of twelve for the manual integration tests passes should be seen within an appropriate context. The manual tests check happy cases, some edge cases, such as permission denied and enforcing desktop lockdown. However, they do not cover hundreds of simultaneous users or fuzzing the API itself or penetration testing the security provided by the Electron application itself. These are all legitimate limitations we have explicitly mentioned in the limitations part of the conclusion.

    In terms of qualitative results, the fact that the integrated solution has indeed decreased the cognitive load mentioned as the issue during the survey of other applications was found as very important in the experiment. If a teacher gives a lab quiz using Invigilo, he or she does not have to master two different systems. Similarly, if a student takes the quiz, he or she does not need two separate room codes/credentials.
  }
\end{justify}

\section{Objective Achieved}

\begin{justify}
  {\fontsize{12}{18}\selectfont
    Table 4.5 maps each of the original objectives stated in Section 1.3 to the evidence available in the current build of the system. The CLO mapping column lists the course learning outcomes that the corresponding objective contributes to.
  }
\end{justify}

\begin{table}[H]
  \caption{Objective and CLO Mapping}
  \label{tab:4.5}
  \centering
  {\fontsize{10}{13}\selectfont
    \begin{tabular}{|p{2in}|p{2in}|p{0.7in}|p{0.6in}|}
      \hline
      \textbf{Objective} & \textbf{Evidence} & \textbf{Status} & \textbf{CLO} \\
      \hline
      Three role architecture with protected endpoints & Admin/Teacher/Student split with role middleware & Achieved & CLO1, CLO2 \\
      \hline
      Desktop fullscreen lock during active exam & Electron main process enforcement verified in TC-11 & Achieved & CLO2, CLO3 \\
      \hline
      Violation detection and reporting & Fullscreen exit, focus loss, blocked shortcut, multiple display, forbidden process detectors & Achieved & CLO2, CLO4 \\
      \hline
      Mixed question type support & MCQ, written and coding share one submission record & Achieved & CLO2 \\
      \hline
      Multi language code execution with timeouts & Node.js, Python and g++ with 8s and 12s limits & Achieved & CLO2, CLO3 \\
      \hline
      Six character room code joining & Generation, uniqueness check and student onboarding flow implemented & Achieved & CLO2 \\
      \hline
      Webcam face detection events & BlazeFace events and snapshot pipeline integrated & Achieved & CLO2, CLO3 \\
      \hline
      Live teacher controls per participant & Freeze, unfreeze and force submit endpoints and UI & Achieved & CLO2, CLO4 \\
      \hline
      Manual evaluation desk with auto recalculation & Implemented and verified in TC-08 & Achieved & CLO2 \\
      \hline
      Persistent audit trail & Submissions, violations, snapshots and evaluations stored in PostgreSQL & Achieved & CLO2, CLO4 \\
      \hline
      Comprehensive documentation & This report and the in repository guides & Achieved & CLO6 \\
      \hline
    \end{tabular}
  }
\end{table}

\section{Conclusion}

\begin{justify}
  {\fontsize{12}{18}\selectfont
    Implementation demonstrated that all the design decisions taken previously were feasible. In fact, the system is built, operates, and performs as predicted in the design part. The quantitative performance metrics are well within the acceptable range for the departmental scale deployment of the system, while the manual testing process is successfully completed from end to end. The limitations, which are observed in terms of adversary resistance, containerisation, and load testing, have all been clearly identified and a future team can take up those areas. Chapter 3 examines the implications of running a system like this in a university environment.
  }
\end{justify}

%==============================================================
% CHAPTER V: Impact Analysis
%==============================================================
\chapter{Societal, Health, Environment, Safety, Ethical, Legal and Cultural Issues}

\section{Intellectual Property Considerations}

\begin{justify}
  {\fontsize{12}{18}\selectfont
    The Invigilo system architecture is based on an open source stack. These include Electron, React, Vite, Node.js, Express, PostgreSQL, TensorFlow.js, BlazeFace, and Monaco Editor, all of which are available under permissive open source licenses like MIT and Apache 2.0. The team included no proprietary library that will create any licensing conflict for the academic and educational purpose stated in this document.

    Any code developed by the team for this project, including the Electron main process security, the room code generator, the manual evaluation station, and the live invigilation software, is the intellectual property of the authors and the Department of Computer Science \& Engineering of KUET, as is customary practice with regards to intellectual property ownership of CSE 3200 class project work. In case this codebase is made public at any point in the future, it can be licensed as an open source library under MIT or similar without any problem related to upstream dependency licenses.

    Question content authored by teachers using the platform, and answer content submitted by students, remains the intellectual property of the respective parties. The system stores this content for the purpose of evaluation and audit, not for commercial reuse. If the platform is later extended to share question banks across departments, an explicit consent and licensing flow should be added before any cross sharing is enabled.
  }
\end{justify}

\section{Ethical Considerations}

\begin{justify}
  {\fontsize{12}{18}\selectfont
    Administering a proctored examination is far from an ethically neutral activity. In terms of technical aspects, the Invigilo application watches the camera, processes system calls, and locks the computer. All of these things constitute measurable invasions of privacy on the part of a user. However, the stance adopted by Invigilo in relation to such privacy invasions is that such invasions are permissible only if they are explicitly delimited, clear to the user, and integral to the integrity of the test.

    The use of the webcam feature in Invigilo is entirely optional and can be turned off by any teacher for any quiz. When a quiz does involve the use of a webcam, it will be made explicitly known to the student before taking the quiz. Only a limited number of snapshots is made of limited size, and events rather than raw video are generated for institutional defense later on.

    The scanner for the prohibited processes operates based on a predefined list of process names. The list has been recorded and can be examined by anyone who accesses the code. There is no machine learning algorithm used to subjectively determine whether a process is suspicious or not. This makes it easy for students to know what they are not supposed to do during the test.

    The controls that have been given to teachers like the freeze or the force submit button have too much power and can easily be misused. However, it is mandatory to record the time when these controls are invoked and by which teacher as a way of having an audit trail. In future, the program will include a prompt message when the force submit is pressed.
  }
\end{justify}

\section{Safety Considerations}

\begin{justify}
  {\fontsize{12}{18}\selectfont
    From a physical safety perspective, the project does not introduce risks of the kind that a hardware engineering project would. There are no high voltage components, no moving parts and no radio emitters. The relevant safety considerations for Invigilo are around mental and emotional safety during a high pressure exam.

    A strict lockdown that triggers many violations in quick succession can put a student under unnecessary stress, especially if the student is not familiar with the tool. Two design decisions try to mitigate this. First, the warning display has a four second timeout so that violation messages do not pile up on the screen. Second, the threshold for automatic forced submission is set at three violations rather than at one, to give the student a margin for genuine accidents such as a brief notification popup from the operating system.

    During a real classroom deployment the teacher should also have the ability to override an auto submit on the spot. This control is available through the evaluation desk by reopening a submission for review, and it is also possible to allow a re attempt for a specific student on a specific exam through the room code flow. These are administrative escape hatches that protect the student from being locked out of an exam by a false positive.
  }
\end{justify}

\section{Legal Considerations}

\begin{justify}
  {\fontsize{12}{18}\selectfont
    From a physical safety perspective, the project does not introduce risks of the kind that a hardware engineering project would. There are no high voltage components, no moving parts and no radio emitters. The relevant safety considerations for Invigilo are around mental and emotional safety during a high pressure exam.

    A strict lockdown that triggers many violations in quick succession can put a student under unnecessary stress, especially if the student is not familiar with the tool. Two design decisions try to mitigate this. First, the warning display has a four second timeout so that violation messages do not pile up on the screen. Second, the threshold for automatic forced submission is set at three violations rather than at one, to give the student a margin for genuine accidents such as a brief notification popup from the operating system.

    Before a production rollout in any university, three legal artefacts should be added. The first is an explicit informed consent flow that the student signs once per semester, explaining what is collected and how it will be used. The second is a written retention policy that states the maximum number of days that snapshots and events are kept. The third is a documented incident response plan in case a snapshot is leaked outside the institution.
  }
\end{justify}

\section{Impact of the Project on Societal, Health, and Cultural Issues}

\begin{justify}
  {\fontsize{12}{18}\selectfont
    The impact of such software in society is largely positive from an educational perspective. If the software is implemented properly, it helps to ensure that exams can be administered fairly, give feedback to students in time, save on paper, and leave behind an audit log that would come in handy during disputes. An unbiased assessment process adds value to the academic certificates issued by the department since this, in turn, influences the perception that employers and other educational institutions have of the department’s graduates.

    Culturally speaking, the software complements rather than undermines the classroom culture already present in the department. The teacher continues to control the test. Using the room code option reflects how an educator in a Bangladeshi lab class conducts an oral announcement of quiz codes to a class comprising thirty students. Similarly, locking the desktop reflects discipline in taking a traditional sit down exam where the focus is entirely on the paper provided. In effect, Invigilo digitalizes an existing culture.

    From the health standpoint, too much screen time and webcam surveillance may result in eye fatigue and mild anxiety. These issues arise from any online exam and are not unique to Invigilo itself. The solution involves procedural measures rather than technological ones, such as limited time of the test, breaks for more extended tests, and instructions in case of feeling sick during the exam.

    Regarding cultural inclusiveness, the interface uses straightforward English and excludes any idiomatic phrases that would confuse a non-native speaker. Future improvements to the software would include the addition of the Bengali language without having to modify the core architecture.
  }
\end{justify}

\section{Impact of Project on the Environment and Sustainability}

\begin{justify}
  {\fontsize{12}{18}\selectfont
    Digital testing in general decreases the use of paper, the printing and photocopying of test questions and answer sheets, and the delivery of test questions and answer sheets physically to the departments. In the case of a department that conducts forty to fifty tests per semester in all sections, there is a significant decrease in paper usage and power spent by photocopiers. Invigilo is a part of this effort by replacing the paper test with a desktop application.

    Conversely, there is an increase in energy expenditure in having to run a computer or laptop continuously with its webcam active throughout the exam process. It uses more energy than scanning a question from a piece of paper. In particular, the webcam recording process ensures that the camera and some cores on the CPU are constantly working for the whole duration of the exam. Moreover, the Electron environment is not one of the most efficient runtimes.

    While the difference in energy consumed is insignificant for a lab quiz with thirty people who take only forty minutes to complete their test, it becomes noticeable for a large campus-wide installation where thousands of students take lengthy tests weekly. A future iteration can refine the snapshot frequency, lower the resolution of the camera to just enough for facial recognition purposes, and switch off the proctoring module altogether in cases where it is unnecessary.

    From a hardware sustainability perspective, the system runs on standard student laptops without requiring any specialised hardware. There is no need to purchase additional sensors, cameras or dongles, which keeps the e waste profile of the deployment close to zero.
  }
\end{justify}

%==============================================================
% CHAPTER VI: Complex Engineering Problems and Activities
%==============================================================
\chapter{Addressing Complex Engineering Problems and Activities}

\section{Complex Engineering Problems Associated with the Current Project}

\begin{justify}
  {\fontsize{12}{18}\selectfont
    A project is considered to address a complex engineering problem when it involves competing constraints, multi domain integration, partial information and risk to stakeholders. Invigilo has all of these characteristics. The most important complex engineering problems addressed during the project are summarised in Table 6.1.
  }
\end{justify}

\begin{table}[H]
  \caption{Complex Engineering Problem Analysis for Invigilo}
  \label{tab:6.1}
  \centering
  {\fontsize{10}{13}\selectfont
    \begin{tabular}{|p{1.4in}|p{1.6in}|p{1.6in}|p{0.9in}|}
      \hline
      \textbf{Complexity Area} & \textbf{Engineering Challenge} & \textbf{Implemented Solution} & \textbf{Residual Risk} \\
      \hline
      Security versus usability & Strong lockdown without breaking the user flow & Electron exam mode lock with bounded warnings and override paths & Privileged users can bypass at OS level \\
      \hline
      Heterogeneous assessment support & Single submission model for MCQ, written and coding & Type tagged answer rows with mixed scoring pipeline & Increased controller complexity \\
      \hline
      Real time governance at lab scale & Teacher needs live information without websocket complexity & Polling based participant and proctoring panels with intervention endpoints & Polling overhead at very large class sizes \\
      \hline
      Proctoring robustness vs privacy & Detect cheating without intrusive data collection & Bounded snapshots, structured events, opt in webcam & False positives, privacy concerns \\
      \hline
      Deterministic fairness & Random order without irreproducibility & Hash seeded shuffle per student per exam & Requires careful key management \\
      \hline
      Safe arbitrary code execution & Run student code in three languages with bounded risk & Spawned processes with strict timeouts and per request directories & Lacks container isolation and fine grained quotas \\
      \hline
      Cross layer correctness & Electron main, renderer and backend must agree on exam state & Single source of truth in PostgreSQL with idempotent transitions & State desync if backend is restarted mid exam \\
      \hline
    \end{tabular}
  }
\end{table}

\begin{justify}
  {\fontsize{12}{18}\selectfont
    Each row in the table represents a real design decision that the team had to defend during the iterative review cycles. The most subtle of these was the security versus usability trade off. An overly aggressive lockdown produces too many false positive violations and frustrates the student. A too relaxed lockdown defeats the purpose of having a desktop client at all. The current configuration, with a violation threshold of three and a four second warning timeout, is the result of several rounds of internal testing and is documented so that another group can adjust it for their own context.
  }
\end{justify}

\section{Complex Engineering Activities Associated with the Current Project}

\begin{justify}
  {\fontsize{12}{18}\selectfont
    Beyond solving discrete problems, the project also involved several complex engineering activities. These are activities that required the team to integrate knowledge from multiple courses, to use industry standard tools, and to coordinate as a team across different layers of the stack.

    \begin{itemize}
        \item \textbf{System Architecture Design:} The three tier architecture, the choice of Electron over a pure web stack, and the decision to keep the security layer in the main process required design reasoning that draws on operating systems, computer networks and software engineering courses simultaneously.

        \item \textbf{Database Schema Engineering:} The seven core tables, the foreign keys, the additive migration strategy and the indexes for common query patterns required applied database design skills.

        \item \textbf{API Contract Design:} Thirty nine endpoints across seven modules were designed with consistent naming, consistent error responses and consistent role checks. Maintaining this consistency across two developers required active discussion and code review.

        \item \textbf{Integration of a Pretrained Model:} BlazeFace integration required understanding the input and output shapes of the model, the timing of inference inside a React component and the size of the snapshots being sent to the backend.

        \item \textbf{Cross Process Communication:} The Electron preload bridge had to expose a narrow set of methods to the renderer without leaking access to the underlying Node.js APIs. This required care to avoid security regressions.

        \item \textbf{Concurrent Process Management:} The coding execution endpoint spawns child processes, applies timeouts, captures output and cleans up temporary files. Doing this safely under load is a non trivial systems programming activity.

        \item \textbf{Real Classroom Validation:} The most underrated complex activity was running the system end to end with classmates as test users and reacting to their feedback. This is what surfaced the keyboard input focus bug in Electron fullscreen mode that led to the modal redesign.
    \end{itemize}

    Together, these activities cover the kind of multi layer, multi tool engineering work that the CSE 3200 course is designed to encourage.
  }
\end{justify}

%==============================================================
% CHAPTER VII: Conclusions
%==============================================================
\chapter{Conclusions}

\section{Summary}

\begin{justify}
  {\fontsize{12}{18}\selectfont
    Invigilo is a self-hosted desktop examination application developed exclusively for the needs of an engineering department in Bangladesh. It comprises a secure main process built with Electron, a React-based renderer engine, a Node.js/Express server backend, and a PostgreSQL database layer. Invigilo accommodates three user roles, three types of questions, and an active teaching interface with freeze and force-submit capabilities. Additionally, it features a webcam-based invigilation pipeline employing BlazeFace detection, a multilingual code evaluation endpoint with tight timeout restrictions, a six-character room code login protocol, and an evaluation desk capable of re-calculating the overall mark dynamically. The latest version of Invigilo reveals thirty-nine API endpoints distributed among seven modules, stores data in seven relational tables, and meets all twelve integration tests conducted manually.
  }
\end{justify}

\section{Limitations}

\begin{justify}
  {\fontsize{12}{18}\selectfont
    The current version of Invigilo has several known limitations that should be made explicit so that future work can address them in a focused way.

    \begin{itemize}
        \item \textbf{No automated test suite:} The verification was done through manual integration runs. There is no unit test layer for the backend controllers and no end to end automation for the renderer.
        \item \textbf{Polling based coordination:} The teacher controls and the participant panel use HTTP polling. For very large classes with hundreds of students this is not the most efficient mechanism and may need to be replaced with websockets.
        \item \textbf{Code execution is process isolated, not container isolated:} The coding endpoint relies on host installed compilers and on operating system process timeouts. A container based sandbox would be safer.
        \item \textbf{Forbidden process scanner is Windows centric:} The list of forbidden process names is tuned for Windows. Cross platform support requires per platform process name lists.
        \item \textbf{Privileged operating system bypasses are not prevented:} A user with administrator rights on the local machine can in principle disable parts of the lockdown. This is an inherent limit of any user space application.
        \item \textbf{Proctoring false positives:} BlazeFace can occasionally classify a head turn as looking away. The teacher has to use judgement when reviewing flagged events.
        \item \textbf{Privacy artefacts are not yet governed by a formal policy:} A written retention policy and a consent template should be added before any production deployment.
        \item \textbf{Single sign on is not integrated:} The system uses its own authentication. Integration with the institutional identity provider is left for future work.
        \item \textbf{Frontend components are large:} The main dashboard files are several thousand lines long and would benefit from decomposition into smaller subcomponents.
    \end{itemize}
  }
\end{justify}

\section{Recommendations and Future Works}

\begin{justify}
  {\fontsize{12}{18}\selectfont
    The natural next steps for extending Invigilo follow directly from the limitations listed above and from the broader needs of a university wide deployment.

    \begin{itemize}
        \item Build an automated test pipeline for the backend API and for the renderer, with continuous integration on every pull request.
        \item Replace the polling based teacher controls with a websocket channel so that freeze and force submit take effect with sub second latency even for large cohorts.
        \item Containerise the coding execution endpoint using Docker with strict CPU, memory and filesystem quotas. This would also make multi language support easier to extend.
        \item Generalise the forbidden process scanner with per platform process name lists for Windows, macOS and Linux.
        \item Add a formal privacy and retention policy with a signed consent flow at the start of every semester.
        \item Integrate with institutional single sign on so that students do not need a separate password for the examination platform.
        \item Decompose the large dashboard components into smaller, individually testable subcomponents.
        \item Add an analytics layer that interprets violation trends across exams and flags possible false positive patterns for teacher review.
        \item Extend the user interface with Bengali language support to make the platform more accessible to students and teachers from a wider range of backgrounds.
        \item Add a mobile companion application for teachers so that they can monitor a live exam without being tied to a desktop.
        \item Explore a question bank import and export format so that teachers across departments can share authored items.
    \end{itemize}

    Even with the limitations acknowledged, the current implementation already demonstrates that a desktop centred, role aware, mixed format examination platform with integrated proctoring is a feasible direction for university level digital assessment in Bangladesh. We hope that future student groups under CSE 3200 will pick the project up from this point and push it closer to a production ready system that the department can deploy at scale.
  }
\end{justify}

%==============================================================
% APPENDIX
%==============================================================
\appendix
\chapter{User Manual and Operational Guide}

\begin{justify}
  {\fontsize{12}{18}\selectfont
    This appendix gives a short operational walkthrough of the Invigilo system for the three user roles.

    \textbf{A.1 System Initialization.} Start the PostgreSQL service and verify that the credential details in \texttt{backend/.env} are accurate. Run \texttt{npm start} to start the system which launches the Electron wrapper, React renderer and Express backend at the same time. For authentication, there are two options: \texttt{Sign In} for accessing all features or \texttt{Enter Exam} for entering the exam with room code.

    \textbf{A.2 Authentication and Mode Access.} The \texttt{Sign In} mode allows the user to authenticate as Admin, Teacher and Student using username/password combination. For students, both email and seven digit roll number can be used in conjunction with the password. The \texttt{Enter Exam} mode requires the room code and student roll number to get into the room.

    \textbf{A.3 Administrator Operations.} After signing in as Admin, dashboard will show statistics for number of teachers, students and exams conducted. Teachers and students can be created through the corresponding dashboard tabs. An import tab is available which uploads a CSV file containing list of students and creates accounts accordingly. Users can be activated/deactivated through the list of users.

    \textbf{A.4 Teacher Operations.} The teacher operations interface is based on the exam list. Creating an exam will lead to the management of the exam interface where the teacher can specify the title, description, exam type, duration, presentation and webcam use for the exam. The question manager allows the teacher to create multiple choice, written and coding questions. The automatically generated room code can be shared among the students. Once the participants are present, the exam will start from the live room interface. Throughout the exam, the participant panel displays the participants' status, violation count, proctoring snapshot and control actions like per participant freeze and force submit. Once the exam finishes, the evaluation desk interface allows the teacher to evaluate written and coding answers.

    \textbf{A.5 Student Operations.} The student can join either via login or room code. The status will automatically update in the waiting room. Once the exam begins, the student responds to the questions as specified by the presentation mode. If it is a coding question, the student chooses the coding language, codes in the provided editor, provides custom input, runs the code and checks the standard output and error message. The student can submit the exam manually or it is automatically submitted when the timer expires. The \texttt{My Results} interface displays the total score and question-wise results.

    \textbf{A.6 Proctoring and Security Behaviour.} While the exam is in progress, the computer screen is fixed in fullscreen always on top. Any violation, such as fullscreen exit, focus loss, and shortcut block will be logged and exposed in the student interface through warnings. If the teacher freezes the participant, then the student will see blocking overlays. If the teacher force-submits the participant, then the force submit overlays will be triggered and the participant will be frozen.

    \textbf{A.7 Common Troubleshooting.} The login failure issue can be sorted out by checking the role selection and credentials, as well as the backend status. Code join failure may have its root cause in incorrect code format and inactive exam state. Missing live updates in the waiting room can occur because of API issues and expired tokens. Failing coding run will happen when there is no proper runtime/compiler installed. Failure of a teacher command due to permission errors will depend on exam ownership and valid teacher token.

    \textbf{A.8 Recommended Operator Checklist.} Prior to taking an exam, create the exam, fill it with questions, set up all required options and test it with one student. While the exam is taking place, check the participant panel, violations, and proctoring panels. After the completion of the exam, make sure that the student sees his/her result.
  }
\end{justify}

%==============================================================
% REFERENCES
%==============================================================
\newpage
\addcontentsline{toc}{chapter}{References}

\begin{center}
  {\fontsize{14}{21}\bfseries REFERENCES}
\end{center}

\blanklines{2}

\begin{spacing}{1.2}
  \setlength{\parskip}{6pt}
  \begin{list}{}{
    \setlength{\leftmargin}{0.31in}
    \setlength{\itemindent}{-0.31in}
    \setlength{\topsep}{0pt}
    \setlength{\itemsep}{6pt}
    \setlength{\parsep}{0pt}
  }
    \fontsize{12}{14}\selectfont

    \item[1.] Electron Foundation, ``Electron Documentation,'' 2026. [Online]. Available: https://www.electronjs.org/docs/latest. [Accessed: Apr. 22, 2026].

    \item[2.] Meta Open Source, ``React Documentation,'' 2026. [Online]. Available: https://react.dev/. [Accessed: Apr. 22, 2026].

    \item[3.] OpenJS Foundation, ``Node.js Documentation,'' 2026. [Online]. Available: https://nodejs.org/en/docs. [Accessed: Apr. 22, 2026].

    \item[4.] StrongLoop and IBM, ``Express Documentation,'' 2026. [Online]. Available: https://expressjs.com/. [Accessed: Apr. 22, 2026].

    \item[5.] PostgreSQL Global Development Group, ``PostgreSQL 16 Documentation,'' 2026. [Online]. Available: https://www.postgresql.org/docs/. [Accessed: Apr. 22, 2026].

    \item[6.] M. Jones, J. Bradley, and N. Sakimura, ``JSON Web Token (JWT),'' RFC 7519, Internet Engineering Task Force, May 2015.

    \item[7.] N. Provos and D. Mazieres, ``A Future Adaptable Password Scheme,'' in Proc. USENIX Annual Technical Conference, 1999, pp. 81--92.

    \item[8.] V. Bazarevsky, F. Kartynnik, A. Vakunov, K. Raveendran, and M. Grundmann, ``BlazeFace: Sub millisecond Neural Face Detection on Mobile GPUs,'' arXiv preprint arXiv:1907.05047, 2019.

    \item[9.] Microsoft, ``Monaco Editor Documentation,'' 2026. [Online]. Available: https://microsoft.github.io/monaco-editor/. [Accessed: Apr. 22, 2026].

    \item[10.] Safe Exam Browser Consortium, ``Safe Exam Browser Documentation,'' 2025. [Online]. Available: https://safeexambrowser.org/. [Accessed: Apr. 22, 2026].

    \item[11.] OWASP Foundation, ``Application Security Verification Standard 4.0.3,'' 2021.

    \item[12.] P. A. Grassi et al., ``Digital Identity Guidelines: Authentication and Lifecycle Management,'' NIST Special Publication 800-63B, National Institute of Standards and Technology, 2017 (with subsequent updates).

    \item[13.] UNESCO, ``Guidance for Generative AI in Education and Research,'' UNESCO Publishing, 2023.

    \item[14.] D. Abrahams, T. Lai, and R. Vasudevan, ``Online Proctoring Practices in Higher Education during and after the Pandemic: A Survey,'' Computers and Education, vol. 195, 2023.

    \item[15.] M. Hussein, M. Yusuf, A. Kabir, and A. Sazali, ``Exam Cheating in Online Assessment: Causes and Mitigation Strategies,'' Education and Information Technologies, vol. 27, no. 5, pp. 6593--6614, 2022.

    \item[16.] A. Krizhevsky, I. Sutskever, and G. E. Hinton, ``ImageNet Classification with Deep Convolutional Neural Networks,'' in Proc. NeurIPS, 2012.

    \item[17.] G. van Rossum and Python Software Foundation, ``The Python Language Reference, Release 3.12,'' 2024.

    \item[18.] B. Stroustrup, ``The C++ Programming Language,'' 4th ed., Addison-Wesley, 2013.

    \item[19.] TensorFlow Authors, ``TensorFlow.js Documentation,'' 2026. [Online]. Available: https://www.tensorflow.org/js. [Accessed: Apr. 22, 2026].

    \item[20.] Government of Bangladesh, ``Digital Security Act, 2018,'' Ministry of Posts, Telecommunications and Information Technology, 2018.

  \end{list}
\end{spacing}

\end{document}
